%  LaTeX support: latex@mdpi.com 
%  For support, please attach all files needed for compiling as well as the log file, and specify your operating system, LaTeX version, and LaTeX editor.

%=================================================================
\documentclass[land,article,submit,pdftex,moreauthors]{Definitions/mdpi} 

%=================================================================

% MDPI internal commands
\firstpage{1} 
\makeatletter 
\setcounter{page}{\@firstpage} 
\makeatother
\pubvolume{1}
\issuenum{1}
\articlenumber{0}
\pubyear{2026}
\copyrightyear{2026}
%\externaleditor{Academic Editor: Firstname Lastname}
\datereceived{16 May 2026} 
\daterevised{14 June 2026} %
\dateaccepted{} 
\datepublished{} 
\hreflink{https://doi.org/} % If needed use \linebreak
%\doinum{}

%=================================================================
% Add packages and commands here. The following packages are loaded in our class file: fontenc, inputenc, calc, indentfirst, fancyhdr, graphicx, epstopdf, lastpage, ifthen, lineno, float, amsmath, setspace, enumitem, mathpazo, booktabs, titlesec, etoolbox, tabto, xcolor, soul, multirow, microtype, tikz, totcount, changepage, attrib, upgreek, cleveref, amsthm, hyphenat, natbib, hyperref, footmisc, url, geometry, newfloat, caption

\graphicspath{{figures/}} % path to folder with figures
\usepackage{subcaption}
\usepackage{listings}
\usepackage{xcolor}
\usepackage{gensymb} % for \textdegree
\usepackage{tabularx}
\newcolumntype{Y}{>{\centering\arraybackslash}X}
\usepackage{amsmath,mathtools,amsthm}
	\setcounter{MaxMatrixCols}{15}
\usepackage{array}
\usepackage{amssymb}
\usepackage{amsfonts}
\usepackage{float}
\usepackage{chngcntr}
\counterwithin*{equation}{section}
\usepackage[super]{nth}
\usepackage{longtable}
\usepackage{tabularx}
\usepackage{multirow}
\usepackage{adjustbox}
\usepackage{ltablex}
\keepXColumns
\usepackage{booktabs}
\usepackage{url}

\definecolor{deepblue}{rgb}{0,0,0.5}
\definecolor{deepred}{rgb}{0.6,0,0}
\definecolor{deepmagenta}{RGB}{195,12,214}
\definecolor{deepgreen}{RGB}{29,122,30}
\definecolor{mintgreen}{RGB}{192,249,192}
\definecolor{codepurple}{RGB}{204, 0, 153}
\definecolor{LightCyan}{rgb}{0.88,1,1}
\definecolor{GainsboroPurple}{RGB}{247,176,247}
\definecolor{LightSlateGrey}{RGB}{124,118,118}
%    
\lstdefinestyle{mystyle}{
    commentstyle=\color{LightSlateGrey},
    keywordstyle=\color{deepgreen}\bfseries,
    emphstyle=\color{deepred},
    numberstyle=\tiny\color{deepblue},
    stringstyle=\color{codepurple},
    basicstyle=\ttfamily\scriptsize,
    breakatwhitespace=false,         
    breaklines=true,                 
    captionpos=b,                    
    keepspaces=true,                 
    numbers=left,                    
    numbersep=5pt,                  
    showspaces=false,                
    showstringspaces=false,
    showtabs=false,                  
    tabsize=2
}
\lstset{style=mystyle}

\numberwithin{equation}{section}


%=================================================================
% Full title of the paper (Capitalized)
\Title{Distributed Deep Learning and Intelligent Soil--Water Analytics in Precision Agriculture: A Comprehensive Review}

% MDPI internal command: Title for citation in the left column
\TitleCitation{Distributed Deep Learning and Intelligent Soil--Water Analytics in Precision Agriculture: A Comprehensive Review}

% Author Orchid ID: enter ID or remove command
\newcommand{\orcidauthorA}{0000-0002-5759-1089} % Polina Add \orcidA{} behind the author's name

% Authors, for the paper (add full first names)
\Author{Polina Lemenkova $^{1,2,3,4}$*\orcidA{}}

% MDPI internal command: Authors, for metadata in PDF
\AuthorNames{Polina Lemenkova}

% MDPI internal command: Authors, for citation in the left column
\AuthorCitation{Lemenkova, P.}

% Affiliations / Addresses
\address{%
$^{1}$ \quad Bureau de Recherches G\'eologiques et Mini\`eres (BRGM). 3 avenue Claude-Guillemin, 45060 Orl\'eans, France.

$^{2}$ \quad LE STUDIUM, Institut d'Etudes Avanc\'ees du Val de Loire. 1 rue Dupanloup 45000, Orl\'eans, France.

$^{3}$ \quad Le laboratoire PRISME, Universit\'e d'Orl\'eans, Institut National des Sciences Appliqu\'ees de Centre-Val de Loire (INSA CVL, UR 4229), Château de la Source, 45100, Orl\'eans, France. 

$^{4}$ \quad La Technopole d'Orl\'eans. 1 avenue du Champ de Mars, 45100 Orl\'eans, France.
}

% Contact information of the corresponding author
\corres{Correspondence: polina.lemenkova@tech-orleans.fr ; Tel.: (+33)0633148601 (P.L.)}

% Abstract (Do not insert blank lines, i.e. \\) 
\abstract{Efficient management of soil--water resources is critical for global food security under intensifying climatic and demographic pressures. This review provides a comprehensive synthesis of artificial intelligence (AI) and distributed deep learning methodologies applied to soil--water interactions in precision agriculture. \textcolor{blue}{The physical and hydraulic foundations of soil--water systems---including water retention, unsaturated flow governed by the Richards equation, and soil degradation processes---are examined and situated within a unified framework of AI-based modelling and decision support.} Classical machine learning (ML) algorithms (random forests, support vector machines, gradient boosting) and deep learning architectures (convolutional neural networks, long short-term memory networks, transformers) are evaluated with respect to their capacity to predict soil moisture dynamics, estimate hydraulic properties, support smart irrigation scheduling, and generate digital soil maps at field-to-regional scales. Distributed training paradigms, federated learning for privacy-preserving multi-farm analytics, and edge AI deployment on low-power IoT hardware are assessed as enabling infrastructures for scalable agricultural intelligence. The review further addresses explainability, uncertainty quantification, and ethical dimensions inherent to AI-driven agricultural systems. Key challenges---including training data scarcity in data-poor regions, model interpretability, integration with physics-based hydrological models, and real-time deployment constraints---are critically discussed. Prospective research directions encompass physics-informed neural networks, foundation models for earth observation, autonomous digital twins of soil--water systems, and federated learning architectures aligned with data sovereignty frameworks. The synthesis underscores AI's transformative potential for sustainable agricultural water management while delineating the technical and sociotechnical barriers that must be resolved to realise this potential at global scale.}

% Keywords
\keyword{Soil-water interactions; machine learning; Deep learning; Distributed learning; Precision agriculture; Smart irrigation; Soil hydrology; Artificial intelligence; Remote sensing; Digital agriculture} 

% The fields PACS, MSC, and JEL may be left empty or commented out if not applicable
%\PACS{07.05.Mh, 07.05.Kf, 92.40.Kf, 92.40.Xx, 89.20.-a, 89.60.-k}
%\MSC{68T07, 68T09, 86A08, 86A32, 62M45, 92B20, 68T05}
%\JEL{Q15, Q25, Q55, C45, O32}

\begin{document}

%%%%%%%%%%%%%%%%%%%%%%%%%%%%%%%%%%%%%%%%%%
%% SECTION 1: INTRODUCTION
%%%%%%%%%%%%%%%%%%%%%%%%%%%%%%%%%%%%%%%%%%

\section{Introduction}

\subsection{Global Water Crisis and Agricultural Sustainability}

Water scarcity has emerged as one of the most pressing and multidimensional challenges confronting global agricultural systems in the twenty-first century. Agriculture accounts for approximately 70\% of total global freshwater withdrawals, making it by far the largest consumer of freshwater resources worldwide \citep{fao2020state}. Concurrently, a confluence of accelerating forces---rapid population growth, urbanisation, changing dietary preferences, and the intensifying effects of anthropogenic climate change---is placing unprecedented pressure on the hydrological systems that sustain food production. Projections by the Intergovernmental Panel on Climate Change (IPCC) indicate that altered precipitation regimes, increased frequency of extreme drought events, and accelerated glacial retreat will substantially reduce water availability across key agricultural regions in Asia, Sub-Saharan Africa, the Mediterranean basin, and parts of the Americas by mid-century \citep{ipcc2022ar6}. In this context, the development of intelligent, data-driven frameworks for efficient soil--water management is no longer merely desirable---it is a scientific and societal imperative.

The challenge of sustainable water management in agriculture is compounded by the heterogeneity of soil systems, the spatial and temporal variability of precipitation, and the complex feedback between land use, vegetation, and the hydrological cycle, Figure \ref{fig01}.

\begin{figure}[H]
    \begin{center}
	\hspace{0pt}\includegraphics[width=14.0cm]{Figure_01.png}
    \end{center}
    \caption{Global overview of water scarcity and agricultural water demand under climate change scenarios. \textcolor{blue}{Software used: QGIS v. 3.44 for mapping; R v. 4.6.1 and ggplot2 package v. 4.0.3 for plotting bar charts and pie charts. Inkscape, v. 1.4.4 for dashboard styling and assembling. Dashboard} source: author.}
    \label{fig01}
\end{figure}

Traditional irrigation practices, which frequently rely on fixed schedules and empirical rules, are ill-suited to the dynamic and location-specific character of soil--water interactions. Globally, it is estimated that between 40\% and 60\% of water applied through irrigation is lost to inefficient distribution, evaporation, deep percolation, and surface runoff \citep{galland2025climate}. Such inefficiencies not only deplete scarce freshwater reserves but also accelerate soil salinisation, waterlogging, and nutrient leaching---degradation processes that undermine long-term agricultural productivity.

\subsection{Importance of Soil--Water Interactions}

At the heart of agricultural water management lies the intimate and dynamic relationship between soil physical properties and water movement. Soil--water interactions govern a broad suite of processes that are fundamental to plant growth and ecosystem function: infiltration, redistribution, evapotranspiration, capillary rise, groundwater recharge, surface runoff generation, and nutrient transport. The capacity of a soil to retain and transmit water is determined by its texture, structure, organic matter content, bulk density, and pore-size distribution---properties that vary considerably across landscapes and that are sensitive to management practices and climatic forcing.

The hydraulic behaviour of soils is inherently nonlinear and spatially heterogeneous. In the unsaturated (vadose) zone, water retention and conductivity functions exhibit complex, history-dependent relationships that make deterministic modelling extremely demanding. The van~Genuchten--Mualem model and the Brooks--Corey parametrisation have been widely adopted to describe water retention curves, yet their parameterisation requires intensive field and laboratory measurement campaigns that are impractical at regional or continental scales \citep{vangenuchten1980closed}. Similarly, processes such as preferential flow through macropores, crack flow in shrink-swell soils, and lateral redistribution in layered profiles introduce further complexity that classical Richards-equation-based simulators struggle to capture adequately without very fine spatial discretisation.

Understanding these dynamics is essential not only for optimising irrigation efficiency but also for predicting the hydrological response of catchments to rainfall events, assessing the risk of contaminant transport to groundwater, and designing effective soil conservation measures. Soil degradation---through salinisation, compaction, erosion, and loss of organic carbon---directly diminishes the water-holding capacity of agricultural soils, amplifying both drought stress during dry periods and runoff-related flooding during intense rainfall events. Climate-driven shifts in precipitation patterns are expected to exacerbate these impacts, particularly in already water-stressed regions, underscoring the urgent need for adaptive soil and water management strategies informed by robust predictive tools \citep{lehmann2020soil}.

\subsection{Emergence of AI and Data-Driven Agriculture}

The past decade has witnessed a technological transformation in the methods available for monitoring, modelling, and managing soil--water systems, Table \ref{tab:evolution_models}. The proliferation of low-cost Internet of Things (IoT) sensor networks, unmanned aerial vehicles (UAVs), and high-resolution satellite remote sensing platforms---including Sentinel-1/2, Landsat-8/9, MODIS, and SMAP---has resulted in an unprecedented abundance of spatially and temporally dense environmental data \citep{ran}. Simultaneously, advances in cloud computing infrastructure, edge computing, and high-performance computing have democratised access to the computational resources required to process and analyse these large, heterogeneous datasets. These twin developments have created fertile conditions for the rapid adoption of artificial intelligence (AI) and \textcolor{blue}{machine learning} (ML) methods across the agricultural and environmental sciences.

%============= TABLE ==================================>
\footnotesize
\begin{longtable}{p{1.4cm} p{1.6cm} p{2.2cm} p{2.9cm} p{2.9cm} p{1.5cm}}
\caption{Evolution of modelling approaches in soil--water systems from empirical and physics-based models to modern artificial intelligence and distributed deep learning frameworks.}
\label{tab:evolution_models} \\
\toprule
\textbf{Era} & \textbf{Modelling Approach} & \textbf{Representative Models} & \textbf{Strengths} & \textbf{Limitations} & \textbf{Key Sources} \\
\midrule
\endfirsthead

\caption[]{Evolution of modelling approaches in soil--water systems (continued).} \\
\toprule
\textbf{Era} & \textbf{Modelling Approach} & \textbf{Representative Models} & \textbf{Strengths} & \textbf{Limitations} & \textbf{Key Sources} \\
\midrule
\endhead

\bottomrule
\endfoot

1950s--1980s &
Empirical and Statistical Models &
Linear regression, empirical infiltration equations, water balance models &
Simple implementation, low computational requirements, easy interpretation, suitable for small datasets &
Limited capability for nonlinear processes, poor generalization under heterogeneous environmental conditions &
\citep{Philip1957, Green1911, Thornthwaite1955} \\ \hline

1980s--2000s &
Physics-Based Hydrological Models &
SWAT, HYDRUS, STANMOD, MODFLOW, Richards equation, Darcy-based models &
Physically interpretable, process-based understanding, strong theoretical foundation, effective for hydrological simulations &
Require extensive calibration, sensitive to parameter uncertainty, computationally expensive for large-scale applications &
\citep{Arnold1998, Simunek, Harbaugh2005, Richards1931} \\ \hline

2000s--2015 &
ML Models &
Support Vector Machines (SVM), Random Forest (RF), Decision Trees, kNN, Artificial Neural Networks (ANN) &
Capable of modeling nonlinear relationships, robust prediction accuracy, reduced dependence on physical assumptions &
Require high-quality datasets, feature engineering often necessary, limited interpretability in some models &
\citep{Breiman2001, Maier2000, Elshorbagy2010} \\ \hline

2015--2020 &
Deep Learning Approaches &
Convolutional Neural Networks (CNN), Recurrent Neural Networks (RNN), Long Short-Term Memory (LSTM) &
Excellent spatiotemporal learning capability, automatic feature extraction, strong performance with remote sensing and time-series data &
Large training datasets required, high computational cost, risk of overfitting, black-box behavior &
\citep{Kratzert2018, Reichstein2019, Shen2018} \\ \hline

2020--Present &
Advanced AI and Distributed Intelligence &
Transformers, Federated Learning, Edge AI, Physics-Informed Neural Networks (PINNs), Distributed Deep Learning &
Real-time analytics, scalability, decentralized learning, privacy preservation, integration with IoT and smart farming systems &
Infrastructure complexity, high energy consumption, limited explainability, challenges in standardization and deployment &
\citep{Karniadakis2021, Lai} \\ \hline

Future Trends &
Hybrid Intelligent Systems &
Digital twins, Explainable AI (XAI), autonomous irrigation systems, climate-integrated AI frameworks &
Adaptive decision-making, improved sustainability, interpretable predictions, integration with climate resilience strategies &
Requires interdisciplinary integration, lack of benchmark datasets, regulatory and ethical concerns &
\citep{Pylianidis2021, Verdouw2021, arrieta2020explainable, Sit2020} \\

\end{longtable}
\normalsize
%============= TABLE ==================================<

AI-driven frameworks offer several capabilities that are difficult or impossible to achieve with conventional physically-based models. Deep learning architectures---including convolutional neural networks (CNNs), long short-term memory (LSTM) networks, and transformer-based attention models---are capable of learning highly complex, nonlinear mappings between input variables (soil properties, meteorological forcing, remote sensing signals) and target outputs (soil moisture content, evapotranspiration rates, crop yield) directly from observational data \citep{Tripathy}. This is achieved without requiring the explicit specification of process equations, making such models particularly attractive for systems where mechanistic understanding remains incomplete or computationally intractable \citep{lecun2015deep}. These models can be trained on historical datasets of arbitrary size, updated in real time as new observations become available, and deployed across diverse agroecological contexts through transfer learning and domain adaptation techniques.

The concept of precision agriculture---tailoring agronomic inputs to the specific needs of individual field zones in space and time---has been substantially advanced by AI-based decision-support systems. Intelligent irrigation controllers that integrate real-time sensor data, numerical weather prediction outputs, and ML-based soil moisture forecasts are now being evaluated in research settings and increasingly in commercial deployments \citep{vij2020iot}. Such systems hold the potential to reduce agricultural water consumption substantially while maintaining or improving crop yields, contributing to the twin sustainability goals of food security and water conservation.

\subsection{Distributed and Federated Deep Learning}

The scale of modern agricultural data introduces both opportunities and challenges for AI-based modelling. Individual farms or monitoring stations may generate millions of sensor readings per year, while satellite platforms produce petabytes of earth observation data globally. Training deep learning models on such volumes of data at centralised facilities is increasingly impractical, both because of the cost of data transmission and storage and because of emerging regulatory constraints on the cross-border transfer of commercially sensitive agricultural data, Table \ref{tab:distributed_ai}. Distributed deep learning---in which model training is parallelised across multiple compute nodes or edge devices---has emerged as a key technological response to these challenges \citep{ben2019demystifying}.

%============= TABLE ==================================>
\setcounter{table}{1}
\footnotesize
\setlength{\tabcolsep}{4pt}
\renewcommand{\arraystretch}{1.25}

\begin{longtable}{p{1.4cm} p{1.4cm} p{3.2cm} p{3.2cm} p{2.4cm} p{1.6cm}}
\caption{Distributed AI frameworks and agricultural applications.}
\label{tab:distributed_ai} \\
\toprule
\textbf{Framework} &
\textbf{Architecture Type} &
\textbf{Key Characteristics} &
\textbf{Agricultural Applications} &
\textbf{Major Challenges} &
\textbf{Key Sources} \\
\midrule
\endfirsthead

\caption[]{Distributed AI frameworks and agricultural applications (continued).} \\
\toprule
\textbf{Framework} &
\textbf{Architecture Type} &
\textbf{Key Characteristics} &
\textbf{Agricultural Applications} &
\textbf{Major Challenges} &
\textbf{Key Sources} \\
\midrule
\endhead

\midrule
\multicolumn{6}{r}{\footnotesize\textit{Continued on next page}} \\
\midrule
\endfoot

\bottomrule
\endlastfoot

Cloud-Based Deep Learning &
Centralized Distributed Computing &
High computational power, scalable storage, parallel processing, large-scale analytics &
Large-scale soil moisture prediction, satellite image analysis, crop monitoring &
Latency, data transfer costs, dependence on internet connectivity &
\citep{Kamilaris2018, Liakos2018} \\ \hline

Edge AI Systems &
Edge Computing &
Real-time processing, low latency, localized intelligence, reduced bandwidth usage &
Smart irrigation control, autonomous sensors, real-time field monitoring &
Limited computational resources and energy constraints &
\citep{Tariq, s25175302} \\ \hline

Federated Learning &
Decentralized Collaborative Learning &
Privacy-preserving distributed training, local data retention, collaborative model optimization &
Multi-farm analytics, privacy-sensitive agricultural data sharing, distributed crop modelling &
Communication overhead and model synchronization complexity &
\citep{Yang3298981, Li2020_FL} \\ \hline

Internet of Things (IoT)-Integrated AI &
Sensor-Based Distributed Systems &
Continuous environmental monitoring, real-time data acquisition, automated analytics &
Precision irrigation, soil moisture sensing, climate monitoring &
Sensor reliability, cybersecurity risks, heterogeneous data integration &
\citep{Paria2025, Khanna2019, Ayaz2019} \\ \hline

Distributed Deep Learning Clusters &
Parallel GPU/Cloud Clusters &
High-performance distributed training, scalable neural network optimization &
Remote sensing analytics, climate prediction, large-scale hydrological simulations &
High infrastructure cost and energy consumption &
\citep{Dean2012, Verbraeken2020} \\ \hline

Fog Computing Frameworks &
Hybrid Cloud--Edge Systems &
Intermediate processing layer between cloud and edge devices, reduced latency &
Real-time agricultural analytics, adaptive irrigation systems &
Complex deployment architecture and maintenance challenges &
\citep{Bonomi2012, s21175922, Guardo} \\ \hline

Blockchain-Integrated AI Systems &
Secure Distributed Intelligence &
Secure decentralized data sharing, transparency, traceability &
Agricultural supply chains, water-use monitoring, smart contracts for irrigation management &
Computational overhead and scalability limitations &
\citep{Ferrag2020, nayak2025blockchain, Casino2019, Reyna2018} \\ \hline

Digital Twin Platforms &
Virtual Simulation Frameworks &
Real-time virtual representation of agricultural systems, predictive optimization &
Soil-water simulation, crop management optimization, climate adaptation &
High data requirements and integration complexity &
\citep{Pylianidis2021, Verdouw2021, ALVES2023135920} \\ \hline

Multi-Agent AI Systems &
Autonomous Distributed Agents &
Collaborative intelligent agents for adaptive decision-making &
Autonomous irrigation scheduling, robotic agriculture, precision resource allocation &
Coordination complexity and communication overhead &
\citep{Wooldridge2009, su17125423, Chlingaryan2018} \\

\end{longtable}
\normalsize
%============= TABLE ==================================<

\subsection{Research Gap and Objectives}

Despite the rapid growth of the literature on AI applications in agriculture, a comprehensive and integrative synthesis of the field remains lacking. Most existing reviews address either ML methods for soil property prediction, AI-based irrigation management, or remote sensing applications in isolation, without situating them within a unified conceptual framework that encompasses distributed computation, explainable AI, uncertainty quantification, and the physical principles governing soil--water systems. Furthermore, the practical implications of privacy-preserving AI architectures and edge-deployed deep learning models for field-scale water management have received limited systematic attention.

This review addresses these gaps by providing a structured synthesis of recent developments across four interrelated domains: 
\begin{enumerate}[label=(\roman*)]
    \item the physical and hydraulic fundamentals of soil--water systems;
    \item AI and deep learning methodologies applicable to soil science and agricultural water management;
    \item distributed and federated learning frameworks for large-scale agricultural analytics; and
    \item explainability, uncertainty, and ethical dimensions of AI deployment in precision agriculture.
\end{enumerate}

The overarching aim is to provide researchers, practitioners, and policymakers with an integrated perspective on the current state of the art and the most promising directions for future research and development.

\textcolor{blue}{This manuscript is organised as follows. 
\begin{enumerate}[label=(\roman*)]
\item Section~1 is the introductory part of the manuscript.
\item Section~2 establishes the physical and hydraulic fundamentals of soil--water systems, covering water retention theory, unsaturated flow governed by the Richards equation, and the principal soil degradation processes relevant to agricultural water management. 
\item Section~3 surveys the AI and deep learning methodologies applicable to soil science and agricultural water management, including classical ML algorithms, deep learning architectures, and physics-informed hybrid models. 
\item Section~4 reviews the data sources underpinning AI-based soil--water modelling, encompassing ground-based sensor networks, satellite and UAV remote sensing platforms, and data fusion challenges. \item Section~5 presents applications of AI in soil--water interactions, covering soil moisture prediction, smart irrigation, digital soil mapping, pedotransfer functions, and crop yield forecasting. 
\item Section~6 provides a comparative analysis of model performance and discusses generalisation and transferability across environments. 
\item Section~7 addresses explainability, uncertainty quantification, and the ethical and socioeconomic dimensions of AI deployment in precision agriculture. 
\item Section~8 critically examines current challenges and limitations.
\item Section~9 delineates prospective research directions and technical horizons. 
\item Section~10 presents the conclusions.
\end{enumerate}}

\subsection{\textcolor{blue}{Review Methodology and PRISMA Compliance}}

\textcolor{blue}{This article is a comprehensive narrative review of the application of artificial intelligence and distributed deep learning to soil--water systems in precision agriculture. It constitutes a systematic meta-analysis where the review methodology is aligned with the PRISMA 2020 guidelines \citep{regulation2016eu} to the extent applicable to comprehensive narrative reviews, in order to ensure transparency, reproducibility of the literature search, and rigorous reporting of the selection process. A PRISMA-adapted flow diagram is provided in Figure~\ref{fig:prisma}.}

\textcolor{blue}{\textbf{Databases and sources.} The primary literature search was conducted across six major peer-reviewed publication platforms and citation indexing databases: ScienceDirect (Elsevier), Wiley Online Library, Taylor \& Francis Online, MDPI Open Access Journals, Scopus, and Web of Science (WoS). These platforms were selected to provide comprehensive coverage of the soil science, hydrology, precision agriculture, and computer science literature. Supplementary records were identified through citation searching of key review papers, examination of reference lists of retrieved articles, and expert recommendations within the author's research network.}

\textcolor{blue}{\textbf{Search terms.} Search queries were constructed by combining domain-specific thematic terms---``soil moisture'', ``soil hydraulic properties'', ``vadose zone'', ``pedotransfer functions'', ``irrigation scheduling'', ``precision agriculture'', ``digital soil mapping'', ``soil water content'', ``evapotranspiration''---with methodological and computational terms---``machine learning'', ``deep learning'', ``convolutional neural network'', ``long short-term memory'', ``transformer'', ``federated learning'', ``physics-informed neural network'', ``explainable AI'', ``uncertainty quantification'', ``distributed learning'', ``edge computing''. Boolean AND/OR operators were used to construct structured queries within each database interface.}

\begin{figure}[H]
\begin{center}
\includegraphics[width=15.0cm]{Figure_PRISMA.pdf}
\end{center}
\caption{\textcolor{blue}{PRISMA 2020 flow diagram (adapted for comprehensive narrative review) documenting the literature identification, screening, eligibility assessment, and inclusion process. Adapted from: Page MJ, et al. BMJ 2021;372:n71. doi:10.1136/bmj.n71. Licensed under CC BY 4.0. Software used for plotting diagram:} \LaTeX \space tikz package. Diagram source: author.}
\label{fig:prisma}
\end{figure}

\textcolor{blue}{\textbf{Temporal scope.} The search covered publications from January 2000 to May 2026. Priority was assigned to publications from 2015 onwards, reflecting the emergence of modern deep learning methods and their adoption in soil and agricultural sciences. Foundational pre-2015 works establishing theoretical frameworks (Richards equation, van~Genuchten model, HYDRUS, MODFLOW, classical ML algorithms) were included irrespective of publication date given their continued relevance to the field.}

\textcolor{blue}{\textbf{Inclusion criteria.} Records were included if they: (i) were published in peer-reviewed journals or peer-reviewed conference proceedings; (ii) addressed at least one of the four thematic domains defined in the review objectives; (iii) were available in English; and (iv) reported sufficient methodological detail to allow assessment of the AI approach and its application context. Grey literature, preprints, and non-peer-reviewed reports were excluded unless they represented canonical technical documentation with no peer-reviewed equivalent (e.g., USGS technical manuals, FAO technical reports, IPCC assessment reports).}

\textcolor{blue}{\textbf{Screening and selection process.} An initial pool of 10,480 database records was identified, supplemented by 205 records from other sources. After deduplication and automated pre-screening for relevance (removing 6,370 records), 4,315 records were retained for title and abstract screening. Full-text retrieval was sought for 1,259 records that passed the screening stage; 35 were not retrievable. Of the 1,224 full-text reports assessed for eligibility, 974 were excluded for the reasons detailed in Figure~\ref{fig:prisma}. The final reference corpus comprises \textbf{250 sources}, grouped thematically into the section structure of this review. The complete selection process is documented in the PRISMA-adapted flow diagram below.}

%%%%%%%%%%%%%%%%%%%%%%%%%%%%%%%%%%%%%%%%%%
%% SECTION 2: FUNDAMENTALS
%%%%%%%%%%%%%%%%%%%%%%%%%%%%%%%%%%%%%%%%%%

\section{\textcolor{blue}{Fundamentals of Soil--Water Systems in Agriculture}}

\subsection{Soil Physical and Hydraulic Properties}

A thorough understanding of soil physical and hydraulic properties is foundational to any quantitative treatment of soil--water dynamics and forms an essential prerequisite for the effective application of AI-based predictive models. The key physical descriptors of a soil include its particle-size distribution (texture), aggregate structure, bulk density ($\rho_b$), total porosity ($\phi$), and organic matter content, Table \ref{tab:soil_hydraulic_properties}. Together, these properties determine the geometry of the pore network through which water and solutes move and within which microbial communities and plant roots develop \citep{hillel2004, Lemenkova202318, HILLEL2003407, Lemenkova202224}.

%============= TABLE ==================================>
\footnotesize
%\setcounter{table}{1}
\setlength{\LTleft}{0pt}
\setlength{\LTright}{0pt}

\begin{longtable}{@{\extracolsep{\fill}} p{0.14\textwidth} p{0.22\textwidth} p{0.26\textwidth} p{0.22\textwidth} p{0.11\textwidth} @{}}
\caption{Major soil hydraulic properties and their agricultural significance.}
\label{tab:soil_hydraulic_properties} \\
\toprule
\textbf{Hydraulic Property} &
\textbf{Definition} &
\textbf{Agricultural Significance} &
\textbf{Factors Affecting Soil Property} &
\textbf{Key Sources} \\
\midrule
\endfirsthead

\caption[]{Major soil hydraulic properties and their agricultural significance (continued).} \\
\toprule
\textbf{Hydraulic Property} &
\textbf{Definition} &
\textbf{Agricultural Significance} &
\textbf{Factors Affecting Soil Property} &
\textbf{Key Sources} \\
\midrule
\endhead

\bottomrule
\endfoot

Porosity &
Fraction of soil volume occupied by pores or voids &
Controls water storage capacity, aeration, root penetration, and microbial activity &
Soil texture, organic matter, compaction, aggregation, tillage practices &
\citep{Barik2020, DEOLIVEIRA2021104814, Rabot2018} \\ \hline

Bulk Density &
Mass of dry soil per unit bulk volume &
Influences root growth, infiltration, and water retention; high bulk density restricts crop productivity &
Soil compaction, machinery traffic, soil texture, moisture content &
\citep{Assouline2011, Chaudhari2013, Keller2011} \\ \hline

Field Capacity &
Amount of water retained in soil after gravitational drainage &
Represents the upper limit of plant-available water and determines irrigation scheduling &
Texture, structure, organic matter, pore-size distribution &
\citep{Amsili2024, Twarakavi2009, Nachabe1998} \\ \hline

Permanent Wilting Point &
Soil moisture level below which plants cannot extract water &
Defines lower threshold of plant-available water and drought stress conditions &
Clay content, salinity, organic matter, soil mineralogy &
\citep{Gamagedara2025, Kirkham2014, Pidgeon1972} \\ \hline

Available Water Capacity &
Difference between field capacity and wilting point &
Determines the quantity of water accessible for crop uptake and irrigation efficiency &
Texture, rooting depth, soil structure, organic matter &
\citep{Amsili2024, ALTOBELLI20201765, Lal2020, Veihmeyer1927} \\ \hline

Hydraulic Conductivity &
Ability of soil to transmit water under saturated or unsaturated conditions &
Controls infiltration, drainage, groundwater recharge, and irrigation performance &
Texture, pore connectivity, compaction, soil moisture conditions &
\citep{Mualem1976, VanGenuchten1980} \\ \hline

Infiltration Rate &
Rate at which water enters the soil surface &
Affects runoff generation, irrigation uniformity, erosion risk, and water conservation &
Surface crusting, vegetation cover, soil texture, tillage practices &
\citep{Assouline2011, Philip1957, Green1911, w10111581} \\ \hline

Soil Water Retention &
Capacity of soil to retain water against gravitational forces &
Critical for drought resistance, irrigation planning, and crop-water availability &
Clay minerals, organic matter, pore-size distribution &
\citep{Gamagedara2025, VanGenuchten1980, Hodnett2002, Vereecken1989} \\ \hline

Matric Potential &
Energy status of water held within soil pores due to capillary and adsorptive forces &
Influences plant water uptake, unsaturated flow, and root-zone dynamics &
Soil texture, salinity, moisture content, pore geometry &
\citep{Gamagedara2025, Richards1931, Dane2002} \\ \hline

Capillary Rise &
Upward movement of water through soil pores due to capillary forces &
Supports crop water availability in shallow groundwater regions and arid environments &
Pore-size distribution, groundwater depth, soil texture &
\citep{hydrology12030055, Gardner1958, ALI2026, app13063374} \\ \hline

Water Holding Capacity &
Maximum amount of water retained by soil after saturation and drainage &
Determines irrigation frequency, crop resilience, and water-use efficiency &
Organic matter, texture, aggregation, compaction &
\citep{Barik2020, flint2018increasing, Hudson1994, Rawls2003} \\

\end{longtable}
\normalsize
%============= TABLE ==================================<

Soil texture---the relative proportions of sand, silt, and clay fractions---exerts a primary control on hydraulic behaviour, Figure \ref{fig02}. Sandy soils are characterised by large inter-particle pores, rapid drainage, low water-holding capacity, and high saturated hydraulic conductivity ($K_s$). Clay-rich soils, by contrast, possess predominantly micropores and mesopores that confer high water retention and low $K_s$, but are also prone to shrinkage-induced cracking upon desiccation and to surface sealing under rainfall impact. Loamy soils, occupying an intermediate textural position, generally offer the most favourable combination of water retention and aeration for crop production. The United States Department of Agriculture (USDA) and the Food and Agriculture Organization (FAO) textural classification schemes provide internationally standardised frameworks for communicating soil texture, and texture-based pedotransfer functions (PTFs) are routinely used to estimate hydraulic parameters from routinely measured soil properties \citep{wosten2001pedotransfer, Lemenkova202302, Lemenkova202314}.

\begin{figure}[H]
    \begin{center}
	\hspace{0pt}\includegraphics[width=14.0cm]{Figure_02.png}
    \end{center}
    \caption{Conceptual diagram scheme of soil--water interaction processes in agricultural systems. \textcolor{blue}{Software: Inkscape v. 1.4.4 for vector elements, GIMP, v. 3.2.4 for raster elements; OpenClipart (openclipart.org) for free SVG clipart of plants, clouds, sun. Infographic} source: author.}
    \label{fig02}
\end{figure}

Several hydraulic properties are of particular practical importance for irrigation management. The field capacity ($\theta_{FC}$) is the volumetric water content retained by a well-drained soil approximately 24--48 hours after saturation, typically corresponding to a matric potential of $-33$~kPa for most medium-textured soils. The permanent wilting point ($\theta_{WP}$) is the water content at which the matric potential falls to approximately $-1500$~kPa, beyond which most crops can no longer extract water from the soil against the combined resistance of soil and root hydraulic pathways. The plant-available water capacity (PAWC), defined as $\theta_{FC} - \theta_{WP}$, quantifies the reservoir of soil water accessible for crop growth between irrigation events and is a critical parameter in irrigation scheduling models.

Hydraulic conductivity describes the ease with which water moves through the soil under a given hydraulic gradient. Its dependence on water content is described by the unsaturated hydraulic conductivity function $K(\theta)$ or, equivalently, $K(h)$, where $h$ is the soil matric potential. The van~Genuchten--Mualem model \citep{vangenuchten1980closed} provides a widely used analytical description:

\begin{equation}
K(h) = K_s \, S_e^{l} \left[1 - \left(1 - S_e^{n/(n-1)}\right)^{(n-1)/n}\right]^2
\end{equation}

\noindent where $S_e = (\theta - \theta_r)/(\theta_s - \theta_r)$ is the effective saturation, $\theta_r$ and $\theta_s$ are the residual and saturated water contents, $n$ and $\alpha$ are shape parameters, and $l$ is a pore-connectivity parameter commonly set to 0.5. Estimating these parameters for field conditions is a significant practical challenge that AI-based PTFs increasingly seek to address.

\subsection{Water Movement in Saturated and Unsaturated Soils}

Water movement through saturated soils is governed by Darcy's law, which relates the volumetric flux density $q$ (m s$^{-1}$) to the hydraulic gradient:

\begin{equation}
q = -K_s \frac{dH}{dl}
\end{equation}

\noindent where $H$ is the total hydraulic head (sum of matric and gravitational heads) and $l$ is the flow path length. Darcy's law provides an adequate description of slow, laminar flow through the saturated zone, and forms the basis of widely used groundwater models such as MODFLOW. In the unsaturated zone, which is of primary relevance for crop water uptake and irrigation management, flow is described by the Richards equation, obtained by combining Darcy's law with the continuity equation:

\begin{equation}
\frac{\partial \theta}{\partial t} = \frac{\partial}{\partial z}\left[K(h)\left(\frac{\partial h}{\partial z} + 1\right)\right] - S
\end{equation}

\noindent where $\theta$ is the volumetric water content, $t$ is time, $z$ is the vertical coordinate (positive upward), $h$ is the matric potential head, $K(h)$ is the unsaturated hydraulic conductivity, and $S$ is a sink term representing root water uptake \citep{richards1931capillary}. The Richards equation is highly nonlinear---because both $\theta$ and $K$ are strongly nonlinear functions of $h$---and its numerical solution is computationally demanding, particularly when the spatial domain is large and the soil profile is heterogeneous. This computational burden has been a significant driver of interest in AI-based surrogate models that can approximate Richards-equation solutions at a fraction of the computational cost.

Soil water redistribution following rainfall or irrigation events involves a succession of processes: infiltration into the surface layer, downward redistribution by gravity and capillary gradients, lateral preferential flow through macropores and biological channels, redistribution driven by osmotic gradients in saline soils, and capillary rise from shallow water tables. Each of these processes operates at characteristic time and length scales, and their interaction creates spatially and temporally complex patterns of soil moisture distribution that challenge both field measurement and numerical simulation.

\subsection{Soil Degradation and Water Stress}

Soil degradation constitutes one of the most significant threats to the long-term sustainability of agricultural water use \cite{Safriel2007}. Globally, it is estimated that more than one-third of agricultural land is moderately to severely degraded, with salinisation, erosion, compaction, and loss of soil organic carbon (SOC) representing the dominant degradation pathways \citep{oldeman1992global}. Each of these processes exerts distinct effects on soil--water dynamics and, consequently, on crop water availability and irrigation requirements.

Soil salinisation---arising from irrigation with saline water, seawater intrusion into coastal aquifers, or evaporative concentration of dissolved salts in poorly drained soils---reduces plant-available water by increasing osmotic stress and suppressing root water uptake. It has been estimated that salinity-affected soils account for approximately 20\% of irrigated agricultural land worldwide and that this proportion is increasing \citep{nachshon2018cropland}. Soil compaction, caused by heavy machinery traffic and tillage, reduces macroporosity and increases soil strength, thereby impeding root penetration, restricting drainage, and increasing surface runoff \citep{Lemenkova202222}. Erosion by water removes fertile topsoil and reduces the depth of the rooting zone, diminishing the soil's water storage capacity.

Water stress exerts a profound influence on crop physiology, suppressing leaf expansion, stomatal conductance, photosynthesis, and ultimately yield formation. The relationship between soil water deficit and crop response is captured by crop-specific stress functions implemented in widely used crop growth models such as AquaCrop, DSSAT, and APSIM. AI-based models have increasingly been used to learn these stress response functions empirically from observed yield data, bypassing the need for manually specified functional forms \citep{van2019machine}.

\subsection{Conventional Soil--Water Modelling Approaches}

Physically-based hydrological and soil--water models have underpinned agricultural water management research for several decades. Widely used frameworks include the Soil and Water Assessment Tool (SWAT) for watershed-scale water balance and non-point source pollution modelling \citep{arnold1998swat}; HYDRUS-1D/2D/3D for variably saturated flow and solute transport in the vadose zone \citep{simunek2008hydrus}; MODFLOW for regional groundwater flow simulation; and DSSAT and AquaCrop for coupled crop--water productivity modelling. These tools have contributed enormously to the understanding of soil--water processes and to the design of irrigation systems and water conservation strategies.

Nevertheless, physically-based models are subject to several important limitations that motivate the complementary use of AI approaches, Figure \ref{fig03}. 

\begin{figure}[H]
    \begin{center}
	\hspace{0pt}\includegraphics[width=14.0cm]{Figure_03.png}
    \end{center}
    \caption{Comparison between physics-based and AI-based soil-water modelling frameworks. \textcolor{blue}{Software used: \LaTeX \space for equations and TexText plugin -- for rendering \LaTeX \space equations; Inkscape v. 1.4.4 for plotting vector elements; TikZ package in \LaTeX \space for plotting neural network diagram (panel B); Material Design Icons for free icons and symbols. Flowchart} source: author.}
    \label{fig03}
\end{figure}

Empirical and statistical modelling approaches---including multiple linear regression, generalised additive models, and geostatistical methods such as kriging---offer lower computational cost and require fewer parameters, but their predictive accuracy is limited in strongly nonlinear systems and their extrapolation behaviour outside the range of calibration data is unreliable. These limitations have collectively catalysed the rapid growth of ML and deep learning as complementary or alternative tools for soil--water modelling. The limitations of physically-based hydrological models include extensive parameterisation requirements, restricted spatial resolution, potential inaccuracies arising from structural assumptions under complex field conditions, and high computational demands that hinder real-time operational applications:

\begin{itemize}
    \item First, physically-based models require extensive parameterisation. SWAT, for example, may involve hundreds of spatially distributed parameters describing soil properties, land use, and channel hydraulics, many of which cannot be directly measured and must be inferred through computationally intensive calibration procedures.
    \item Second, the spatial resolution of these models is constrained by the availability of input data and computational resources. Representing sub-field heterogeneity in regional or global applications is generally infeasible.    
    \item Third, the model structural assumptions (e.g., the Richards equation and the Green--Ampt infiltration model) may not hold under all conditions encountered in practice, particularly in degraded soils with complex pore geometries or in the presence of non-equilibrium flow phenomena.    
    \item Fourth, the computational cost of running physically-based models at scales relevant to real-time operational management is prohibitive without significant simplification.
\end{itemize}

Hybrid model architectures, which integrate the process knowledge embedded in physically-based models with the pattern-recognition capabilities of ML, are increasingly recognised as a promising path forward. Physics-informed neural networks (PINNs), for example, incorporate differential equation constraints directly into the neural network loss function, ensuring that the learned solution remains consistent with governing physical laws even in data-sparse regions \citep{raissi2019physics}. Similarly, ML models have been used to replace computationally expensive process model components (e.g., the solution of the Richards equation) with fast, accurate emulators that can be integrated into real-time decision systems.

%%%%%%%%%%%%%%%%%%%%%%%%%%%%%%%%%%%%%%%%%%
%% SECTION 3: AI IN SOIL & WATER SCIENCES
%%%%%%%%%%%%%%%%%%%%%%%%%%%%%%%%%%%%%%%%%%

\section{\textcolor{blue}{Artificial Intelligence in Soil and Water Sciences}}

\subsection{Overview of AI Techniques and Conceptual Framework}

Artificial intelligence encompasses a broad spectrum of computational approaches that enable machines to perform tasks that would conventionally require human expertise: pattern recognition, prediction, classification, optimisation, and adaptive decision-making \cite{Fan2018}. Within the context of soil and water sciences, AI methods can be broadly categorised into classical ML algorithms (which rely primarily on handcrafted feature engineering) \cite{Amani}, deep learning architectures (which learn hierarchical feature representations directly from raw data) \cite{Attri}, and hybrid or physics-guided models that integrate AI with domain knowledge \citep{jordan2015machine}. Each category offers distinct advantages and is suited to different problem types, data modalities, and application scales.

A unifying characteristic of all AI approaches is their reliance on historical data for model training: the model parameters (weights, split thresholds, kernel coefficients, etc.) are adjusted to minimise a loss function measuring the discrepancy between model predictions and observed target variables \citep{}. This data-driven paradigm contrasts with the a priori specification of process equations in physically-based models, and its predictive accuracy depends critically on the volume, quality, and representativeness of the training data for classification \cite{Tso2007}. In domains such as soil science where training data may be sparse, geographically biased, or collected with heterogeneous measurement protocols, this dependence introduces significant challenges that have motivated the development of transfer learning, data augmentation, and physics-informed regularisation strategies.

The taxonomy of AI applications in soil and water sciences is increasingly sophisticated, Figure \ref{fig04}. At the input end, AI models ingest data from a diverse range of sources: ground-based sensors measuring soil moisture, electrical conductivity, temperature, and matric potential; meteorological stations providing precipitation, air temperature, humidity, wind speed, and solar radiation records; multispectral, hyperspectral, and synthetic aperture radar (SAR) satellite imagery; UAV-based RGB, multispectral, and thermal imagery; and agronomic records of crop variety, sowing date, fertiliser application, and yield. 

\begin{figure}[H]
    \begin{center}
	\hspace{0pt}\includegraphics[width=14.0cm]{Figure_04.png}
    \end{center}
    \caption{Taxonomy of AI methods applied in soil-water and agricultural systems. \textcolor{blue}{Software used: Inkscape v. 1.4.4 for vector elements; GIMP v. 3.2.4 for raster elements and gradient fill; Vecteezy free tier for schematic icons. Flowchart} source: author.}
    \label{fig04}
\end{figure}

The heterogeneity of these data sources---in spatial and temporal resolution, measurement precision, and data format---poses significant pre-processing and data fusion challenges that must be addressed before AI models can be effectively trained and deployed.

\subsection{Classical ML Algorithms}

Table \ref{tab:ml_algorithms} provides a comprehensive comparative analysis of various ML algorithms utilized in soil--water modelling, evaluating their core architectures, data requirements, and predictive performance across diverse agro-climatic conditions.

%============= TABLE ==================================>
\footnotesize
\setlength{\tabcolsep}{3pt}
\renewcommand{\arraystretch}{1.3}

\begin{longtable}{p{1.4cm} p{1.3cm} p{2.6cm} p{2.6cm} p{2.6cm} p{2.3cm}}
\caption{Comparison of ML algorithms in soil--water modelling.}
\label{tab:ml_algorithms} \\
\toprule
\textbf{Algorithm} &
\textbf{Category} &
\textbf{Strengths} &
\textbf{Limitations} &
\textbf{Typical Applications} &
\textbf{Key Citations} \\
\midrule
\endfirsthead

\caption[]{Comparison of ML algorithms in soil--water modelling (continued).} \\
\toprule
\textbf{Algorithm} &
\textbf{Category} &
\textbf{Strengths} &
\textbf{Limitations} &
\textbf{Typical Applications} &
\textbf{Key Citations} \\
\midrule
\endhead

\midrule
\multicolumn{6}{r}{\footnotesize\textit{Continued on next page}} \\
\midrule
\endfoot

\bottomrule
\endlastfoot

Linear Regression &
Supervised Learning &
Simple implementation, computational efficiency, high interpretability &
Limited ability to capture nonlinear relationships and complex interactions &
Water balance estimation, evapotranspiration prediction, trend analysis &
\citep{Montgomery2021, Allen1998, Gaoetal2017, Konukcu, SCOTT2021108251} \\ \hline

Decision Trees &
Supervised Learning &
Easy interpretation, handles nonlinear relationships, low preprocessing requirements &
Prone to overfitting and instability with noisy datasets &
Soil classification, irrigation decision support, land suitability assessment &
\citep{Quinlan1986, Safavian1991, Padarian2020} \\ \hline

Random Forest (RF) &
Ensemble Learning &
High predictive accuracy, robust against overfitting, effective for high-dimensional datasets &
Reduced interpretability and increased computational cost for large datasets &
Digital soil mapping, soil moisture prediction, crop yield forecasting &
\citep{Breiman2001, Hengl2017, Carranza2021, Liaw2002, Biau2016} \\ \hline

Support Vector Machines (SVM) &
Supervised Learning &
Strong generalization capability, effective with limited datasets, suitable for nonlinear modelling &
Sensitive to parameter selection and kernel choice, computationally intensive for large datasets &
Salinity prediction, irrigation scheduling, soil property estimation &
\citep{Gill2006, Asefa, Behzad, Karandish2016} \\ \hline

k-Nearest Neighbors (kNN) &
Instance-Based Learning &
Simple implementation, nonparametric nature, adaptable to local data patterns &
Sensitive to noise and data scaling, computationally expensive for large datasets &
Soil texture classification, local hydrological prediction &
\citep{Cover1967, Altman1992, VELOSO2022e00569, Elshorbagy2010, Fix1989} \\ \hline

Artificial Neural Networks (ANN) &
Deep Learning &
Excellent nonlinear modelling capability, adaptive learning, robust predictive performance &
Requires large datasets and extensive training, limited interpretability &
Hydraulic conductivity estimation, rainfall-runoff modelling, evapotranspiration forecasting &
\citep{Hornik1989, Maier2000, Govindaraju2000, ASCE2000, Dawson2001} \\ \hline

Convolutional Neural Networks (CNN) &
Deep Learning &
Strong spatial feature extraction capability, highly effective for image analysis and remote sensing &
High computational requirements and large data dependency &
Remote sensing analysis, digital soil mapping, crop-water stress detection &
\citep{LeCun2015, Padarian2019, ai5020033, Krizhevsky2017, Ng2019} \\ \hline

Recurrent Neural Networks (RNN) &
Deep Learning &
Efficient sequence and temporal dependency modelling &
Gradient vanishing problems and limited long-term memory retention &
Time-series soil moisture prediction, hydrological forecasting &
\citep{Jordan1986, Elman1990, Coulibaly2001, hydrology12080207, Medsker2001} \\ \hline

Long Short-Term Memory (LSTM) &
Deep Learning &
Excellent long-term temporal learning capability, effective for sequential environmental datasets &
Computationally intensive and complex hyperparameter tuning &
Soil moisture forecasting, rainfall prediction, irrigation demand estimation &
\citep{Hochreiter1997, Kratzert2018, Fang2017, Zhang2019_LSTM, Aetal} \\ \hline

Gradient Boosting Machines (GBM) &
Ensemble Learning &
High predictive accuracy, efficient handling of heterogeneous datasets &
Sensitive to parameter tuning and potential overfitting &
Crop yield prediction, groundwater quality assessment &
\citep{Subudhi, Friedman2001, Natekin2013, Elith2008} \\ \hline

Extreme Gradient Boosting (XGBoost) &
Ensemble Learning &
Fast training, high scalability, superior predictive performance &
Complex optimization and reduced interpretability &
Precision agriculture, smart irrigation systems, hydrological modelling &
\citep{Chen2016, rs17132129, Ghiat, Benassi} \\ \hline

Transformer Models &
Deep Learning &
Superior spatiotemporal learning capability, parallel processing efficiency, scalable architectures &
Very high computational demand and large training data requirements &
Climate-smart agriculture, large-scale hydrological forecasting, environmental analytics &
\citep{Devlin2018, Wu2021_Trans, Nguyen} \\

\end{longtable}
\normalsize
%============= TABLE ==================================<

\subsubsection{Support Vector Machines}

Support Vector Machines (SVMs) represent a class of kernel-based supervised learning algorithms originally developed by \citet{cortes1995support} for binary classification and subsequently extended to regression tasks (support vector regression, SVR). The core principle of SVM is to identify the hyperplane in a high-dimensional feature space that maximises the margin between classes (for classification) or that contains the largest proportion of training points within a specified $\varepsilon$-tube around the regression surface (for SVR).

For binary classification, the optimal separating hyperplane is obtained by solving:

\begin{equation}
\min_{\mathbf{w},b} \frac{1}{2}\|\mathbf{w}\|^2
\quad
\text{subject to }
y_i(\mathbf{w}^\top \mathbf{x}_i+b)\geq 1,
\end{equation}

where $\mathbf{w}$ is the weight vector, $b$ is the bias term, and $y_i$ denotes the class labels.

The decision function of kernel-based SVM is expressed as:

\begin{equation}
f(\mathbf{x})=\sum_{i=1}^{n}\alpha_i y_i K(\mathbf{x}_i,\mathbf{x})+b,
\end{equation}

where $\alpha_i$ are Lagrange multipliers and $K(\mathbf{x}_i,\mathbf{x})$ is the kernel function.

The ability of SVMs to operate effectively in high-dimensional feature spaces---through the kernel trick, which implicitly maps inputs to high-dimensional representations without explicit computation---makes them particularly well-suited to problems involving spectral remote sensing data, where the number of spectral bands may greatly exceed the number of training samples.

In soil science, SVMs have been successfully applied to soil moisture prediction from meteorological and remote sensing inputs, classification of soil types from multispectral imagery, estimation of soil organic carbon content from hyperspectral reflectance spectra, and prediction of soil salinity distribution from electrical conductivity survey data \citep{viscarra2010predictive, Lemenkova202411}. SVMs are attractive for these applications because their regularisation framework mitigates overfitting even when training sample sizes are limited---a chronic problem in soil survey databases---and because their kernel hyperparameters can be efficiently optimised using cross-validation. However, SVMs scale poorly with training set size (their computational complexity is at least $\mathcal{O}(n^2)$ in the number of training samples), limiting their application to large national or global soil databases without approximation methods.

\subsubsection{Random Forests}

Random Forest (RF) is an ensemble learning algorithm based on the aggregation of multiple decision trees trained on bootstrap samples of the training data and evaluated on randomly selected subsets of predictor variables at each node split \citep{breiman2001random}. By averaging the predictions of a large number of individually weak learners, RF achieves substantially lower variance and higher generalisation accuracy than individual decision trees, while retaining their ability to model complex nonlinear interactions and to handle mixed-type predictor variables without pre-processing.

The RF prediction is computed as the average of predictions from all trees:

\begin{equation}
\hat{y}=\frac{1}{T}\sum_{t=1}^{T} h_t(\mathbf{x}),
\end{equation}

where $T$ is the number of trees and $h_t(\mathbf{x})$ is the prediction of the $t$-th tree.

Feature importance in RF can be estimated using permutation importance:

\begin{equation}
FI_j=\frac{1}{T}\sum_{t=1}^{T}
\left(
\text{Err}_{t,j}^{\text{perm}}-\text{Err}_t
\right),
\end{equation}

where $\text{Err}_{t,j}^{\text{perm}}$ denotes the prediction error after permuting feature $j$.

RF has become one of the most widely used algorithms in digital soil mapping (DSM), driven by its strong empirical performance, its built-in feature importance metric (based on the mean decrease in impurity or permutation accuracy), and its robustness to irrelevant predictor variables and moderate amounts of missing data. In the global DSM project SoilGrids, RF models were trained on a global compilation of soil profile observations (over 150,000 profiles) combined with a comprehensive stack of environmental covariates (climate, terrain, vegetation indices, parent material, land use) derived from remote sensing and global datasets, producing predictions of soil properties at 250~m resolution for all soil layers down to 2~m depth \citep{hengl2017soilgrids250m}. Similar approaches have been applied at national and regional scales to map soil organic carbon, pH, clay content, bulk density, and available water capacity with high spatial resolution \citep{Lemenkova202502, Lemenkova202501}.

Beyond DSM, RF has been applied to soil moisture prediction from combinations of in situ sensor data, satellite-derived land surface temperature and vegetation indices, and meteorological variables; to the estimation of soil hydraulic conductivity from texture, bulk density, and organic matter measurements; and to the identification of soil contamination hotspots from geochemical survey data. The interpretability of RF through variable importance analysis is an important advantage in scientific applications where understanding the drivers of model predictions is as important as the predictions themselves.

\subsubsection{Gradient Boosting Methods and XGBoost}

Gradient boosting is an ensemble technique in which learners are added sequentially, with each new learner fitting the residuals of the ensemble assembled so far \citep{friedman2001greedy}. Unlike RF, which builds trees independently in parallel, gradient boosting constructs trees in a stage-wise fashion, explicitly minimising a differentiable loss function.

The additive boosting model is expressed as:

\begin{equation}
F_m(x)=F_{m-1}(x)+\gamma_m h_m(x),
\end{equation}

where $F_m(x)$ is the boosted model at iteration $m$, $h_m(x)$ is the newly added weak learner, and $\gamma_m$ is the learning rate.

In XGBoost, the objective function combines the training loss and regularisation term:

\begin{equation}
\mathcal{L}=
\sum_{i=1}^{n}
l(y_i,\hat{y}_i)
+
\sum_{k=1}^{K}\Omega(f_k),
\end{equation}

where $l(y_i,\hat{y}_i)$ is the loss function and $\Omega(f_k)$ controls model complexity.

The extreme gradient boosting algorithm (XGBoost) \citep{chen2016xgboost} and its successors LightGBM and CatBoost have emerged as the state-of-the-art among tabular ML algorithms, consistently achieving top performance in data science competitions across a wide range of prediction tasks.

In soil and water science applications, XGBoost and gradient boosting have demonstrated superior performance compared with RF and SVM in several benchmarking studies on soil hydraulic property prediction, irrigation demand forecasting, and soil nutrient mapping. Key advantages include the ability to handle missing predictor values natively, efficient implementation with parallel computation, and built-in regularisation terms that control model complexity. Gradient boosting methods are increasingly combined with SHAP (SHapley Additive exPlanations) values \citep{lundberg2017unified} to provide feature attribution scores that quantify the contribution of each input variable to individual predictions, substantially improving model interpretability without sacrificing predictive accuracy \citep{Lemenkova202518, Lemenkova202505}.

\textcolor{blue}{A critical comparative assessment of classical ML algorithms reveals that their relative performance is highly conditioned on the nature of the prediction task, the volume and quality of training data, and the soil--climatic context of deployment. Random Forest consistently demonstrates the strongest generalisation performance for digital soil mapping and soil property prediction tasks, particularly when training datasets are moderately sized (hundreds to low thousands of observations) and predictor sets include strongly correlated remote sensing covariates, owing to its inherent feature bagging mechanism that mitigates multicollinearity effects} \citep{hengl2017soilgrids250m}. 

\textcolor{blue}{XGBoost and gradient boosting variants outperform RF in dynamic process prediction tasks---such as irrigation demand forecasting and evapotranspiration estimation---where the sequential residual-fitting mechanism provides superior capture of nonlinear temporal dependencies in tabular meteorological records. SVMs offer a statistically principled framework for small-sample problems (fewer than approximately 500 labelled observations) through structural risk minimisation, but their $\mathcal{O}(n^2)$--$\mathcal{O}(n^3)$ training complexity renders them computationally prohibitive for continental-scale soil databases. A key limitation shared across all classical ML approaches---and often understated in the literature---is their inability to provide physically consistent predictions outside the training distribution: they do not encode the conservation laws or differential equation constraints that physically-based models enforce, making their extrapolative reliability under novel climate forcing conditions fundamentally uncertain without physics-informed regularisation \citep{willard2022integrating}.}

\subsection{Deep Learning Architectures}
Representative deep learning architectures for spatiotemporal soil moisture prediction are compared in terms of their approaches to spatial feature learning and temporal dependency modeling. Differences among CNN–LSTM, ConvLSTM, and graph-based temporal frameworks in data representation, network connectivity, and prediction mechanisms are illustrated, Figure \ref{fig05}.

\subsubsection{Convolutional Neural Networks (CNN)}

Convolutional Neural Networks (CNNs) are the dominant deep learning architecture for tasks involving spatially structured data, particularly images \citep{lecun1998gradient}. CNNs exploit the spatial locality of features in images through the use of local receptive fields, weight sharing across spatial positions, and hierarchical feature extraction through successive convolutional, normalisation, and pooling operations. These design principles confer powerful translation invariance and computational efficiency relative to fully connected networks, enabling CNNs to learn rich spatial representations from large image datasets.

\begin{figure}[H]
    \begin{center}
	\hspace{0pt}\includegraphics[width=14.0cm]{Figure_05.png}
    \end{center}
    \caption{Deep learning architectures for spatiotemporal soil moisture prediction. \textcolor{blue}{Software used: Python's Matplotlib library ($mpl\_toolkits.mplot3d$ toolkit for 3D plotting; draw.io (diagrams.net) for plotting LSTM chain diagram. Node-edge graphs with coloured nodes were plotted using Python, NetworkX and Matplotlib. Mathematical notations were rendered in \LaTeX . Final assembly of three panels in layout: Inkscape v. 1.4.4. Flowchart} source: author.}
    \label{fig05}
\end{figure}

In the soil and water sciences, CNNs have found widespread application in the analysis of satellite and UAV imagery for soil classification, soil moisture estimation, crop stress detection, and irrigation mapping. SAR-based soil moisture retrieval has been substantially improved by deep CNN architectures that exploit the spatial context of radar backscatter signals and can account for topographic, vegetation, and surface roughness effects that confound classical retrieval algorithms \citep{hajj2017soil, Lemenkova202402}. Hyperspectral CNNs---designed to process image cubes with hundreds of spectral channels---have achieved state-of-the-art accuracy in soil organic carbon mapping from airborne and proximal sensing data by simultaneously exploiting spectral and spatial information. Transfer learning from large-scale pretrained CNN models (e.g., those trained on ImageNet) to domain-specific soil science tasks has been shown to dramatically reduce the amount of labelled training data required for fine-tuning, facilitating the development of high-accuracy models even in data-scarce agricultural regions.

\subsubsection{Long Short-Term Memory (LSTM) and Recurrent Neural Networks (RNN)}

Many critical soil--water processes are inherently temporal in character: soil moisture evolves in response to sequences of rainfall, evapotranspiration, and irrigation events; streamflow is generated by the cumulative effect of antecedent soil moisture conditions and precipitation over preceding hours to days; crop growth and water demand follow phenological trajectories that unfold over the growing season. Modelling these temporal dependencies requires architectures capable of representing dynamic system state over extended time horizons.

Recurrent Neural Networks (RNNs) address this requirement by maintaining a hidden state vector that is updated at each time step as a function of the current input and the previous hidden state. However, standard RNNs are notoriously difficult to train on long sequences due to the vanishing and exploding gradient problem---the tendency for backpropagated gradients to diminish or amplify exponentially through many time steps \citep{bengio1994learning}. Long Short-Term Memory (LSTM) networks \citep{hochreiter1997long} overcome this limitation through the introduction of a gated memory cell architecture comprising input, forget, and output gates that selectively regulate the flow of information into, within, and out of the cell. This architecture enables LSTMs to learn long-range dependencies spanning hundreds to thousands of time steps, making them exceptionally well-suited to hydrological and soil moisture time-series modelling.

LSTM networks have achieved state-of-the-art performance in rainfall-runoff modelling, outperforming conceptual hydrological models calibrated on thousands of catchments in large-sample benchmarking studies \citep{kratzert2018rainfall}. In soil moisture prediction, LSTMs trained on combinations of meteorological forcing and in situ measurements have demonstrated superior temporal generalisation relative to shallow ML models, capturing soil moisture dynamics across wet and dry periods, freeze-thaw cycles, and inter-annual variability. Bidirectional LSTM (BiLSTM) architectures, which process sequences in both forward and reverse directions, and encoder-decoder LSTMs for multi-step ahead forecasting, represent further refinements that extend the applicability of recurrent architectures to complex temporal prediction tasks in precision agriculture.

\subsubsection{Artificial Neural Networks}

Artificial Neural Networks (ANNs) are computational systems inspired by the organisation of the biological neural network, comprising interconnected processing units (neurons) arranged in layers and trained by backpropagating gradients of a loss function through the network \citep{goodfellow2016deep}. In their shallow multi-layer perceptron (MLP) form---with one or two hidden layers---ANNs have been applied to soil science problems since the early 1990s, initially for estimation of soil hydraulic properties from texture and bulk density \citep{pachepsky1996developing} and subsequently for a wide range of prediction tasks including soil moisture content, evapotranspiration, crop water demand, and water quality indices.

The output of a neuron in an ANN is computed as:

\begin{equation}
a_j = \phi \left( \sum_{i=1}^{n} w_{ij} x_i + b_j \right),
\end{equation}

where $x_i$ are the input features, $w_{ij}$ are the connection weights, $b_j$ is the bias term, and $\phi(\cdot)$ is the activation function (e.g., sigmoid, ReLU, or Leaky ReLU).

The training objective is typically formulated through minimisation of a loss function using backpropagation and gradient descent:

\begin{equation}
\theta^{(t+1)} =
\theta^{(t)}
-
\eta
\frac{\partial \mathcal{L}}{\partial \theta},
\end{equation}

where $\theta$ represents the network parameters (weights and biases), $\eta$ is the learning rate, and $\mathcal{L}$ is the loss function.

Although shallow ANNs demonstrated useful predictive skill in many of these applications, their universal approximation properties are only fully realised with sufficient depth and width, and their performance tends to degrade relative to ensemble tree methods when training data are limited. Deep feedforward networks, with many hidden layers and nonlinear activation functions (ReLU, Leaky ReLU, sigmoid), can learn highly complex mappings, but are prone to overfitting without appropriate regularisation (dropout, weight decay, batch normalisation) and require larger training datasets.

A commonly used regularised loss function for deep ANNs is:

\begin{equation}
\mathcal{L}_{\text{reg}}
=
\mathcal{L}
+
\lambda
\sum_{k} w_k^2,
\end{equation}

where $\lambda$ is the regularisation coefficient controlling weight decay.

Recent work has demonstrated the utility of deep ANNs combined with ensemble techniques for estimating soil water retention curves and hydraulic conductivity functions from basic soil descriptors across diverse global soil datasets \citep{Lemenkova202403, Lemenkova202410}.

\subsubsection{Transformer-Based Architectures and Attention Mechanisms}

Transformer architectures, introduced by \citet{Vaswani2017} for natural language processing, have rapidly emerged as the dominant paradigm for sequence modelling across a broad range of domains, including geoscience and remote sensing. The key innovation of the transformer is the self-attention mechanism, which computes a weighted average of all positions in the input sequence for each output position, with attention weights determined by the compatibility between query and key vectors derived from the input.

The scaled dot-product attention mechanism is defined as:

\begin{equation}
\text{Attention}(Q,K,V)
=
\text{softmax}
\left(
\frac{QK^\top}{\sqrt{d_k}}
\right)V,
\end{equation}

where $Q$, $K$, and $V$ denote the query, key, and value matrices, respectively, and $d_k$ is the dimensionality of the key vectors.

The self-attention weights for each token are computed as:

\begin{equation}
\alpha_{ij}
=
\frac{
\exp \left(
\frac{\mathbf{q}_i^\top \mathbf{k}_j}{\sqrt{d_k}}
\right)
}{
\sum_{m=1}^{n}
\exp \left(
\frac{\mathbf{q}_i^\top \mathbf{k}_m}{\sqrt{d_k}}
\right)
},
\end{equation}

where $\alpha_{ij}$ represents the attention weight assigned by position $i$ to position $j$.

Unlike RNNs and LSTMs, transformers process all elements of the input sequence simultaneously (enabling parallel computation during training) and can capture long-range dependencies without the information bottleneck of a fixed-size hidden state.

Multi-head attention further extends this mechanism by allowing the model to jointly attend to information from multiple representation subspaces:

\begin{equation}
\text{MultiHead}(Q,K,V)
=
\text{Concat}
(\text{head}_1,\dots,\text{head}_h)W^O,
\end{equation}

where each attention head is computed independently as:

\begin{equation}
\text{head}_i
=
\text{Attention}
(QW_i^Q,KW_i^K,VW_i^V).
\end{equation}

In soil and water science, transformer-based models have been applied to time-series prediction of soil moisture, streamflow, groundwater levels, and crop yields, consistently achieving competitive or superior performance relative to LSTM baselines, particularly on long sequences and in multi-variable forecasting tasks. Spatial transformers and vision transformers (ViT) \citep{dosovitskiy2021image} have been adapted for soil classification and moisture mapping from satellite imagery, leveraging the attention mechanism to selectively weight spatially distant pixels with high relevance to the prediction target. Hybrid spatio-temporal transformer architectures---combining spatial attention across the image dimension with temporal attention across the time series dimension---represent the current frontier for AI-based monitoring of dynamic soil--water states from dense time series of satellite observations.

\subsection{Distributed and Federated Deep Learning in Agriculture}

\subsubsection{Distributed Training Frameworks}

Training large deep learning models on agricultural datasets of realistic scale---encompassing multiple seasons, multiple farms, and multiple sensor modalities---demands computational resources that exceed the capacity of a single machine. Distributed deep learning frameworks distribute both the data and the computational workload across multiple processors or nodes, enabling training to complete in hours rather than weeks \citep{ben2019demystifying}. Two principal paradigms are employed: data parallelism, in which each worker processes a distinct shard of the training data and the computed gradients are aggregated (by summation or averaging) before the model weights are updated; and model parallelism, in which the model itself is partitioned across workers, with each worker responsible for the forward and backward passes through a subset of layers, Figure \ref{fig06}.

\begin{figure}[H]
    \begin{center}
	\hspace{0pt}\includegraphics[width=14.0cm]{Figure_06.png}
    \end{center}
    \caption{Architecture of distributed deep learning and federated learning in smart agriculture. \textcolor{blue}{Software used: draw.io built-in Network library for plotting server racks (data centre pics); OpenClipart for icons (UAV/drone, weather station, soil sensor, tractor, market); Noun Project free tier for plotting icons of Smart Agriculture Application panel. Assembly: Inkscape v. 1.4.4. Flowchart} source: author.}
    \label{fig06}
\end{figure}

Large-scale distributed training has enabled the development of foundation models for earth observation---deep learning models pretrained on vast corpora of satellite imagery spanning diverse land cover types, seasons, and geographic regions---that can be fine-tuned with relatively small amounts of task-specific data for downstream applications including soil property mapping, irrigation classification, and crop type delineation. Projects such as the NASA-IBM Prithvi model and ESA's EuroSAT benchmark have demonstrated the value of this approach for accelerating the development of high-accuracy agricultural monitoring systems.

\subsubsection{Federated Learning for Privacy-Preserving Agricultural AI}

\textcolor{blue}{Federated learning, a specific paradigm of distributed AI, allows ML models to be trained collaboratively across a network of decentralised participants---individual farms, regional data hubs, or national agricultural agencies---without requiring the raw data to leave the local node. Only model gradients or parameter updates are communicated to a central aggregator, which combines them to produce a globally improved model. This architecture preserves data privacy and sovereignty while enabling the collective intelligence of a distributed sensor network to be leveraged for model training. In the context of soil--water management, federated learning is particularly promising for aggregating heterogeneous field data from geographically diverse agricultural regions, enabling the development of robust models that generalise across a wide range of soil types and climatic conditions.}

Federated learning (FL) was introduced by \citet{mcmahan2017communication} as a framework for training ML models collaboratively across a network of distributed clients while keeping each client's data local. In the canonical FL algorithm (FedAvg), each communication round proceeds as follows: the central server distributes the current global model parameters to a selected subset of clients; each client performs a fixed number of gradient descent steps on its local data; the updated local parameters are transmitted back to the server; and the server aggregates the local updates (typically by weighted averaging proportional to local dataset size) to produce an improved global model.

The local model update on client $k$ is performed using gradient descent:

\begin{equation}
\theta_k^{(t+1)}
=
\theta_k^{(t)}
-
\eta
\nabla \mathcal{L}_k(\theta_k^{(t)}),
\end{equation}

where $\theta_k$ denotes the local model parameters for client $k$, $\eta$ is the learning rate, and $\mathcal{L}_k$ is the local loss function computed on the client dataset.

The global aggregation step in the FedAvg algorithm is defined as:

\begin{equation}
\theta^{(t+1)}
=
\sum_{k=1}^{K}
\frac{n_k}{n}
\theta_k^{(t+1)},
\end{equation}

where $K$ is the number of participating clients, $n_k$ is the number of samples held by client $k$, and

\begin{equation}
n=\sum_{k=1}^{K} n_k
\end{equation}

is the total number of training samples across all clients.

This process is iterated until the global model converges. The federated optimisation framework substantially reduces the need for centralised data collection while preserving data privacy and enabling collaborative learning across distributed agricultural, environmental, and geospatial datasets.

The application of FL to agricultural AI is motivated by several practical considerations. Agricultural sensor data from individual farms may constitute commercially sensitive information---soil fertility maps, yield records, and irrigation logs can reveal proprietary agronomic practices that farm operators are reluctant to share with competitors or data aggregators. Regulatory frameworks such as the European Union General Data Protection Regulation (GDPR) and national agricultural data governance policies impose legal constraints on the transfer and centralised storage of personal and commercially sensitive data \citep{regulation2016eu}. FL provides a technically and legally robust mechanism for harnessing the predictive power of large distributed agricultural datasets while respecting these constraints.

Challenges specific to FL in agricultural contexts include extreme heterogeneity in the statistical distribution of local data (non-IID data), driven by variation in soil types, climate zones, crop species, and management practices across participating farms; communication bandwidth constraints in rural areas with limited connectivity; and the potential for malicious or malfunctioning clients to degrade the global model through byzantine gradient attacks or hardware failures. Differential privacy mechanisms \citep{dwork2014algorithmic}, secure aggregation protocols, and robust aggregation algorithms that down-weight anomalous gradient contributions have been proposed to address these challenges.

\subsubsection{Edge AI and Internet of Things (IoT) Integration}

The deployment of AI inference at the network edge---on IoT devices, field-deployed microcontrollers, or gateway servers located at the farm level---enables real-time, autonomous decision-making for irrigation control, pest detection, and soil condition monitoring without dependence on reliable cloud connectivity \citep{vij2020iot}. Edge AI is particularly important in smallholder agricultural contexts across developing regions where internet infrastructure is unreliable and the cost of cloud computing is prohibitive. 

\textcolor{blue}{The algorithmic families described in the preceding subsections---classical ML, deep learning, and physics-informed hybrid architectures---underpin a broad spectrum of soil--water applications that are treated in detail in Section~5. These include point-to-regional-scale soil moisture prediction (Section~5.1), AI-driven smart irrigation scheduling and water use optimisation (Section~5.2), estimation of soil hydraulic properties through deep learning pedotransfer functions (Section~5.3), salinity, pollution, and soil remediation (Section~5.4.), digital soil mapping and spatial land-use analytics (Section~5.5), crop yield prediction and precision agriculture (Section~5.6), and climate change adaptation (Section~5.7). An illustrative overview of the AI workflow for smart irrigation and water management is provided in Figure~\ref{fig07}.}

\begin{figure}[H]
    \begin{center}
	\hspace{0pt}\includegraphics[width=14.0cm]{Figure_07.png}
    \end{center}
    \caption{\textcolor{blue}{Conceptual system architecture diagram of} AI workflow for smart irrigation and water management systems. \textcolor{blue}{Software and tools used: GIMP, v. 3.2.4; Phosphor Icons for Circular icon at top (Agri-Planner / Farm Manager); draw.io container shape for plotting panels; Inkscape v. 1.4.4 for circular engine icons and assembling the layout}. Diagram source: author.}
    \label{fig07}
\end{figure}

Compact deep learning models, optimised for deployment on resource-constrained hardware through techniques such as weight quantisation (reducing the numerical precision of model parameters from 32-bit to 8-bit or 4-bit floating point), structured pruning (removing low-importance connections or neurons), and knowledge distillation (training a small student model to mimic the predictions of a large teacher model), can achieve accuracy approaching that of large cloud-deployed models while running on microcontrollers with kilobytes of RAM and milliwatts of power.

Integrated soil--water IoT systems typically combine capacitance or time-domain reflectometry (TDR) soil moisture sensors, temperature and rainfall probes, and plant water status sensors (dendrometers, stem psychrometers) with edge AI processors running lightweight LSTM or transformer models fine-tuned for the specific farm and soil context. These systems close the feedback loop between sensing and actuation, enabling precision variable-rate irrigation that responds to real-time soil moisture dynamics at sub-field resolution.

%%%%%%%%%%%%%%%%%%%%%%%%%%%%%%%%%%%%%%%%%%
%% SECTION 4: DATA SOURCES FOR AI-BASED
%% SOIL–WATER MODELLING
%%%%%%%%%%%%%%%%%%%%%%%%%%%%%%%%%%%%%%%%%%

\section{\textcolor{blue}{Data Sources for AI-Based Soil--Water Modelling}}

\textcolor{blue}{The performance of any AI or deep learning model is fundamentally bounded by the quality, volume, and spatial and temporal representativeness of its training data. In soil--water science, this dependency is particularly consequential because target variables---soil moisture content, hydraulic conductivity, solute concentration, evapotranspiration---are inherently continuous in space and time yet are sampled at discrete, often sparse locations and intervals. Bridging this observational gap requires the integration of data streams from a diverse ensemble of sensing technologies, each with characteristic spatial and temporal resolution, measurement uncertainty, and cost profile. The principal data sources reviewed in this section were identified through the literature search described in Section~1.3, with inclusion prioritising studies reporting quantitative validation of AI models against independently collected field or satellite observations.}

% ----------------------------------------------------------
\subsection{Ground-Based Sensor Networks}
% ----------------------------------------------------------

In situ sensors remain the primary source of direct, high-frequency soil--water measurements and provide the ground-truth observations against which remote sensing retrievals and model predictions are validated. The principal measurement technologies include time-domain reflectometry (TDR), frequency-domain reflectometry (FDR), and capacitance probes for volumetric soil moisture content; tension-cup tensiometers and dielectric matric potential sensors for soil water potential; electrical resistivity tomography (ERT) for spatially continuous subsurface moisture imaging; and electromagnetic induction (EMI) instruments for bulk electrical conductivity surveys of soil salinity \citep{vereecken2014soil}. Thermometric sensors, weighing lysimeters, and eddy-covariance flux towers complement these measurements with continuous records of soil temperature, actual evapotranspiration, and surface energy balance components.

The emergence of low-cost IoT-enabled sensor platforms has dramatically expanded the density and geographic reach of soil moisture networks over the past decade. National and continental networks---including the USDA Soil Climate Analysis Network (SCAN), the International Soil Moisture Network (ISMN), the OzNet network in Australia, and the Cosmic-Ray Neutron Sensing (CRNS) network---now maintain archives of multi-decadal, multi-depth soil moisture records that are publicly accessible and increasingly exploited for ML model training \citep{dorigo2021ismn}. The ISMN alone aggregates data from over 2,600 stations across 71 networks worldwide, providing a globally distributed training corpus for soil moisture prediction models. Wireless sensor networks (WSNs) deployed at the farm scale can resolve sub-field spatial variability at metre-scale resolution, generating high-frequency data streams that feed directly into real-time edge AI inference pipelines \citep{vij2020iot}. However, sensor drift, fouling, cable damage, and power outages introduce systematic gaps and outliers that must be detected and corrected through automated quality-control algorithms before the data can be ingested by ML models \citep{gruber2020}.

% ----------------------------------------------------------
\subsection{Remote Sensing and Satellite Data}
% ----------------------------------------------------------

Satellite remote sensing has transformed the spatial coverage and temporal frequency of soil--water observations, providing synoptic, repeatable measurements over large areas that are inaccessible to ground-based sensor networks. The electromagnetic spectrum from visible and near-infrared to passive microwave and active radar wavelengths encodes information about soil moisture, surface temperature, vegetation water content, and land surface properties, which can be retrieved through empirical algorithms, physically-based forward models, or increasingly, deep learning-based inversion schemes \citep{petropoulos2015}.

Passive microwave missions---particularly the NASA Soil Moisture Active Passive (SMAP) satellite launched in 2015 and the ESA Soil Moisture and Ocean Salinity (SMOS) mission---provide daily global surface soil moisture retrievals at coarse spatial resolution (36~km for SMAP L-band radiometer). SMAP data have been extensively used as training targets for deep learning downscaling models that leverage high-resolution optical and SAR auxiliary data to sharpen retrievals to 1--10~km resolution \citep{fang2019deep}. The Copernicus Sentinel constellation has substantially expanded the repertoire of operationally available remote sensing datasets for agricultural monitoring: Sentinel-1 C-band SAR enables all-weather soil moisture retrieval and crop mapping at 10~m resolution with a 6-day revisit time \citep{hajj2017soil}, while Sentinel-2 provides 10--20~m multispectral imagery with a 5-day revisit that is widely used for vegetation index time-series extraction, crop type classification, and surface energy balance estimation \citep{Lemenkova202405}.

Thermal infrared remote sensing from Landsat-8/9 TIRS, MODIS MOD11 land surface temperature products, and the ECOSTRESS instrument aboard the International Space Station provides estimates of land surface temperature (LST) that, in conjunction with vegetation indices, enable the retrieval of evapotranspiration through the surface energy balance equation and data-driven models such as the Surface Energy Balance Algorithm for Land (SEBAL) and ML-based variants \citep{allen1998crop, song2022deep, Lemenkova202221}. Hyperspectral missions---including the PRISMA satellite launched by the Italian Space Agency in 2019 and the upcoming NASA Surface Biology and Geology (SBG) mission---extend spectral resolution to hundreds of narrow bands, enabling the detection of soil organic carbon, clay mineral composition, iron oxide content, and other geochemical properties through spectral mixture analysis and convolutional neural network regression \citep{hong2021multimodal}.

The integration of observations across satellite platforms at different spatial and temporal scales---a process known as spatiotemporal fusion---is an active research area. Algorithms based on dictionary learning, sparse representation, and encoder-decoder deep learning architectures have been developed to fuse coarse-resolution, high-temporal-frequency data (e.g., MODIS daily reflectance) with fine-resolution, low-temporal-frequency data (e.g., Landsat bi-weekly imagery) to produce synthetic time series at Landsat resolution and MODIS revisit frequency \citep{gao2017blending}. Such fused products substantially increase the density of high-resolution observations available for training and validating AI-based soil--water models in cloudy tropical and temperate regions.

% ----------------------------------------------------------
\subsection{UAV and Drone-Based Monitoring}
% ----------------------------------------------------------

Unmanned aerial vehicles (UAVs) have emerged as a powerful and flexible complement to satellite remote sensing for high-resolution, on-demand monitoring of soil--water dynamics and crop water status at the field scale \citep{mulla2013twenty}. Equipped with RGB, multispectral, hyperspectral, thermal infrared, and LiDAR sensors, UAVs can acquire centimetre-resolution imagery at user-defined acquisition times and revisit frequencies, unconstrained by satellite overpass schedules or cloud cover limitations. This flexibility makes them particularly valuable for capturing transient soil moisture patterns following rainfall or irrigation events, detecting early-stage crop water stress before it manifests in visible symptoms, and generating high-density training datasets for ML models that cannot be obtained from coarser satellite observations alone.

Thermal UAV surveys, in which canopy temperature is retrieved from shortwave infrared cameras, enable the estimation of crop water stress through the Crop Water Stress Index (CWSI) \citep{jones2004use} and serve as inputs to precision variable-rate irrigation controllers. Multispectral UAV data have been used to estimate leaf area index (LAI), fractional vegetation cover, and normalised difference vegetation index (NDVI) at sub-metre resolution, which are key covariates in evapotranspiration models and soil moisture downscaling algorithms. Recent advances in structure-from-motion (SfM) photogrammetry applied to UAV imagery enable the generation of high-resolution digital elevation models (DEMs) and canopy height models that support topographic wetness index computation and precision drainage management \citep{yao2019uav}. The primary constraints on UAV deployment are regulatory restrictions on flight altitude and beyond-visual-line-of-sight operation, battery endurance limiting individual flight coverage, and the substantial processing burden associated with handling very large point clouds and image mosaics.

% ----------------------------------------------------------
\subsection{Big Data Integration and Data Fusion}
% ----------------------------------------------------------

Modern AI-based soil--water modelling draws on an increasingly broad and heterogeneous portfolio of data streams that span multiple observational platforms, spatial scales, thematic domains, and temporal resolutions. Effective exploitation of this data richness requires robust data fusion pipelines that harmonise inputs with different coordinate systems, temporal sampling frequencies, missing data patterns, and measurement uncertainties into a coherent, model-ready feature space \citep{reichstein2019deep}. Cloud computing platforms---including Google Earth Engine, Microsoft Planetary Computer, and the AWS Open Data Registry---have dramatically lowered the barrier to large-scale geospatial data processing by providing petabyte-scale archives of satellite imagery and derived products alongside scalable compute infrastructure, enabling researchers to train and evaluate AI models on global-scale datasets without maintaining local HPC infrastructure \citep{gorelick2017google}.

Data-driven approaches to data fusion, including multimodal deep learning architectures that process inputs from heterogeneous sensor modalities in parallel feature extraction streams before fusing representations in a joint embedding space, have demonstrated superior performance to traditional pixel-based or layer-stacking approaches in agricultural mapping and soil property prediction tasks \citep{hong2021multimodal}. Graph neural networks (GNNs) offer a natural framework for representing the spatial relationships between sensor locations, field parcels, and catchments as graphs, enabling the propagation of information across the observation network in a spatially aware manner \citep{xu2021graph}. The integration of agronomic management records---crop variety, sowing and harvest dates, tillage operations, fertiliser and pesticide applications, and irrigation logs---with geophysical and remote sensing datasets within unified ML frameworks represents a frontier of data fusion that promises substantial improvements in the precision and interpretability of soil--water models \citep{liakos2018machine}.

% ----------------------------------------------------------
\subsection{Data Challenges}
% ----------------------------------------------------------

Despite the expanding volume and diversity of available data, several persistent challenges limit the effectiveness of AI-based soil--water modelling in practice. The most fundamental is class imbalance and spatial sampling bias: existing soil profile databases are disproportionately concentrated in North America, Western Europe, and Australia, leaving vast areas of Africa, Central Asia, and South America severely underrepresented \citep{hengl2017soilgrids250m}. ML models trained on spatially biased datasets exhibit degraded prediction accuracy in underrepresented regions and may propagate systematic errors into global soil property maps and water balance assessments.

Temporal data gaps, arising from sensor failures, cloud cover, satellite overpass gaps, and seasonal inaccessibility of field sites, introduce missingness patterns that are frequently non-random and that pose challenges for time-series ML architectures that require complete input sequences \citep{gruber2020}. Imputation methods based on temporal interpolation, matrix completion, or generative adversarial networks (GANs) have been explored, but their effectiveness depends critically on the nature and extent of the missingness. Label noise---arising from measurement errors in reference soil samples, inconsistent sampling protocols between laboratories and surveys, and the inherent spatial support mismatch between point observations and gridded model predictions---is a further source of degraded training signal that requires explicit treatment through robust loss functions or uncertainty-aware training procedures.

Data sharing and standardisation barriers also constrain the development of large, globally representative training corpora. Harmonising soil profile data collected under different national standards, with different analytical methods and variable nomenclature, requires substantial data curation effort. International initiatives such as ISRIC World Soil Information, the Global Soil Biodiversity Initiative, and the FAO Global Soil Partnership are working to address these challenges through the development of open data standards and federated data infrastructure, but significant interoperability gaps remain.

\section{\textcolor{blue}{Applications of AI in Soil--Water Interactions}}

\subsection{Soil Moisture Prediction}

Soil moisture prediction constitutes one of the most extensively studied applications of AI in soil--water science, driven by the central role of root-zone water content in irrigation scheduling, drought monitoring, and hydrological forecasting \cite{hess2024}. AI-based models have been developed across a hierarchy of spatial and temporal scales, from point-scale time-series prediction at individual sensor stations to regional-scale spatiotemporal mapping using satellite-derived inputs \citep{fang2019deep}.

At the point scale, LSTM networks have demonstrated exceptional capacity to learn the complex nonlinear dynamics of soil moisture evolution in response to precipitation, evapotranspiration, and irrigation inputs \cite{Datta}. \textcolor{blue}{In temporally blocked cross-validated evaluations conducted on multi-year records from point-scale stations distributed across temperate and semi-arid agricultural regions of North America and Europe---drawn from the International Soil Moisture Network (ISMN) and the USDA Soil Climate Analysis Network (SCAN)---LSTM models have achieved root mean square errors (RMSE) in the range of 0.02--0.04~m$^3$~m$^{-3}$ for surface and root-zone volumetric water content at depths of 5--50~cm. These accuracies are competitive with, and in several cases superior to, those of calibrated HYDRUS-1D simulations driven by identical meteorological forcing at the same stations, with the critical distinction that the LSTM models require no prior estimation of van~Genuchten--Mualem hydraulic parameters from laboratory or pedotransfer function data \citep{kratzert2018rainfall}. It should be noted, however, that these results apply specifically to the point scale under the temperate and semi-arid climatic conditions represented in the ISMN and SCAN networks; generalisation to tropical soils, organic soils, or freeze--thaw-affected profiles has not yet been systematically established.} Bidirectional LSTM (BiLSTM) and encoder-decoder architectures extend the predictive horizon to multi-day and multi-week soil moisture forecasts, enabling proactive rather than reactive irrigation management.

At the field scale, hybrid CNN--LSTM models that simultaneously exploit the spatial structure of remotely sensed inputs and their temporal evolution have substantially improved prediction accuracy. These models integrate surface reflectance, land surface temperature, and vegetation indices from Sentinel-2 and MODIS with in situ sensor records to generate spatially distributed soil moisture estimates at sub-field resolution. Transformer architectures incorporating self-attention across both the spatial and temporal dimensions are increasingly competitive with LSTM-based approaches, particularly for long-sequence prediction tasks where the attention mechanism can selectively weight distant historical states.

At the regional scale, deep learning downscaling algorithms have been developed to sharpen coarse-resolution SMAP passive microwave retrievals (36~km footprint) to field-scale resolutions of 1--10~km by incorporating high-resolution SAR backscatter, optical vegetation indices, terrain attributes, and soil texture maps as auxiliary covariates in an encoder-decoder framework \citep{fang2019deep}. Physics-informed neural networks (PINNs) that embed the Richards equation as a soft constraint in the loss function provide a principled mechanism for ensuring that soil moisture predictions remain consistent with mass conservation and hydraulic conductivity relationships, even in data-sparse regions where purely data-driven models tend to extrapolate unreliably \citep{karniadakis2021physics}.

Data assimilation frameworks that combine Kalman filter-based state estimation with AI-based observation operators represent a further frontier, enabling the systematic incorporation of heterogeneous satellite and sensor observations into soil moisture prediction systems with rigorously quantified uncertainty bounds.

\subsection{Smart Irrigation and Water Management}

Intelligent irrigation systems use ML algorithms and sensor feedback to optimize irrigation timing and water allocation, improving water-use efficiency and crop productivity.

Operational smart irrigation systems integrate AI-based soil moisture prediction with real-time IoT sensor telemetry, numerical weather prediction outputs, and crop growth stage information to compute irrigation recommendations that minimise water application while avoiding root-zone moisture deficits that induce yield-limiting crop water stress \citep{vij2020iot}. Reinforcement learning (RL) provides a particularly natural framework for the sequential decision-making structure of irrigation scheduling: an agent observes the current state of the soil--crop--atmosphere system and selects irrigation actions to maximise a cumulative reward encoding crop water productivity and input cost efficiency. \textcolor{blue}{Deep Q-networks (DQN) and actor-critic algorithms (Proximal Policy Optimisation, Soft Actor-Critic) trained within calibrated AquaCrop and DSSAT crop model environments---used as physics-based digital surrogates for agent--environment interaction---have demonstrated irrigation water savings of 20--35\% relative to fixed-schedule and soil-moisture-threshold-based control strategies that are representative of current commercial practice. These savings were reported across simulation experiments covering maize, soybean, and wheat under humid temperate and semi-arid agroclimatic conditions in the United States and the Mediterranean basin, evaluated at the field scale (1--10~ha management units) over growing seasons of 90--180~days. The absence of statistically significant yield penalty under these conditions was established by comparing simulated end-of-season biomass and harvestable yield between RL-controlled and baseline scenarios within the respective crop model environments; it should be emphasised that these outcomes are contingent on the fidelity of the simulation surrogate to real field conditions, and that transferability to heterogeneous field soils, imperfect sensors, and uncertain weather forecasts has not yet been fully validated under operational conditions.}

Model predictive control (MPC) frameworks informed by LSTM-based soil moisture forecasts constitute an alternative AI-driven approach: the controller solves a constrained optimisation problem at each decision step, selecting irrigation amounts that minimise predicted moisture deficit over a rolling forecast horizon subject to water availability and application rate constraints. The integration of probabilistic forecast uncertainty into MPC cost functions---penalising decisions that carry high risk of soil moisture falling below the wilting point---substantially improves the robustness of irrigation recommendations under uncertain weather conditions.

%============= TABLE ==================================>
% Set the table to span the exact width of the page
\footnotesize
% Ensure the table snaps perfectly to the left and right margins of your text layout
\setlength{\LTleft}{0pt}
\setlength{\LTright}{0pt}

\begin{longtable}{p{2.4cm} p{2.3cm} p{2.6cm} p{2.6cm} p{2.6cm} p{2.3cm}}
\caption{AI techniques for smart irrigation and water conservation.}
\label{tab:ai_irrigation} \\
\toprule
\textbf{AI \& Algorithmic Category} & 
\textbf{Specific AI / ML Model} & 
\textbf{Data Source / Inputs} & 
\textbf{Primary Role in Water Conservation} & 
\textbf{Water Saving Metric} \\
\midrule
\endfirsthead

\caption[]{AI techniques for smart irrigation and water conservation (continued).} \\
\toprule
\textbf{AI \& Algorithmic Category} & 
\textbf{Specific AI / ML Model} & 
\textbf{Data Source / Inputs} & 
\textbf{Primary Role in Water Conservation} & 
\textbf{Water Saving Metric} \\
\midrule
\endhead

\midrule
\multicolumn{6}{r}{\footnotesize\textit{Continued on next page}} \\
\midrule
\endfoot

\bottomrule
\endlastfoot

Supervised Learning (Regression \& Prediction) & 
Random Forest (RF), Support Vector Regression (SVR) & 
IoT telemetry (Soil moisture, ambient temp, air humidity, UV index) & 
Predicts the precise localized Surface Soil Moisture (SSM) to prevent over-watering or plant stress. & 
Reduces baseline manual irrigation water use by 30--40\% \citep{Oguzturk2025}. \\
\midrule
Deep Learning \& Sequence Modeling & 
Long Short-Term Memory (LSTM), Transformers & 
Historical agrometeorological time-series, regional weather forecasts & 
Forecasts Daily Reference Evapotranspiration ($ET_0$) and crop water depletion rates days in advance \citep{Zhang2025}. & 
Minimizes evaporative runoff, allowing highly adaptive preemptive irrigation scheduling. \\ 
\midrule
Computer Vision \& Image Processing & 
Convolutional Neural Networks (CNNs), U-Net & 
Multi-spectral satellite imagery (Sentinel/Landsat), UAV/Drone aerial imagery \citep{Yparraguirre2026}. & 
Identifies vegetative indices (NDVI, NDWI) to segment crop water stress zones across large geographic fields. & 
Prevents blanket-field flooding by enabling zone-specific drip irrigation. \\ 
\midrule
Meta-Heuristic Optimization & 
Genetic Algorithms (GA), Particle Swarm Optimization (PSO) & 
Multi-objective constraints (Available water reserves, electricity costs, crop growth stage) \citep{Paria2025}. & 
Runs algorithmic optimization loops to discover the mathematically most efficient water distribution scheduling schedule. & 
Increases overall resource efficiency by 30--50\% across diverse agro-climatic zones \citep{Hamdouni2026}. \\ 
\midrule
Edge \& Intelligent Control & 
Fuzzy Logic Controllers (FLC) & 
Real-time sensor micro-data (pH levels, salinity, soil matric potential) \citep{Ali2025}. & 
Acts as the automated rules-engine that triggers, throttles, or shuts off motorized irrigation valves dynamically without human intervention. & 
Cuts water wastage from system lags; optimizes nutrient-water (``fertigation'') ratios \citep{Oguzturk2025}. \\ 
\midrule
Anomaly Detection \& Diagnostics & 
Autoencoders, Isolation Forests & 
Flow rate meters, acoustic sensors, distributed pressure gauges \citep{w18080919}. & 
Monitors water distribution networks to detect sub-surface pipeline leaks or anomalous pressure drops instantly. & 
Prevents structural water loss and localized over-saturation from damaged plumbing infrastructure. \\
\bottomrule
\end{longtable}

%============= TABLE ==================================<
\normalsize

\subsection{Estimation of Soil Hydraulic Properties}

Deep learning approaches are increasingly used to estimate hydraulic conductivity, infiltration rates, and water retention curves from soil texture and environmental datasets.

Pedotransfer functions (PTFs) are the primary tool for estimating soil hydraulic properties---saturated hydraulic conductivity ($K_s$), water retention curve parameters ($\alpha$, $n$, $\theta_r$, $\theta_s$), and unsaturated conductivity function parameters---from routinely measured basic soil descriptors such as texture fractions, bulk density, and organic matter content \citep{wosten2001pedotransfer}. Classical regression-based PTFs, fitted to relatively small and geographically constrained calibration datasets, exhibit limited transferability across soil types and geographic regions. Deep learning PTFs trained on large globally compiled databases---including UNSODA, HYDRUSbase, and ISRIC World Soil Information archives---have substantially outperformed classical regression PTFs in independent validation studies, particularly for soils with atypical textural compositions or unusual structural properties \citep{zhang2018deep}.

Convolutional neural networks operating on proximal soil sensing data---including visible--near-infrared (Vis-NIR) and mid-infrared (MIR) spectra, X-ray computed tomography (CT) images of soil pore architecture, and electrical resistivity tomography profiles---have achieved state-of-the-art accuracy in predicting $K_s$ and water retention parameters at the sample scale. Graph neural networks (GNNs) that represent the soil pore network as a graph structure and propagate hydraulic information across connected pore bodies and throats offer a physically grounded approach to learning conductivity from pore-scale CT imaging data, bypassing the need for manual extraction of geometric pore parameters. Physics-informed neural network PTFs that constrain predictions to satisfy the van~Genuchten--Mualem model structure---ensuring that estimated water retention and conductivity functions are thermodynamically consistent and monotonic---provide improved extrapolation reliability relative to unconstrained black-box deep learning models \citep{karniadakis2021physics}. Transfer learning from globally pretrained PTF models to locally fine-tuned variants using small regional calibration datasets represents a practically important strategy for deploying high-accuracy hydraulic property estimation in data-sparse agricultural regions.

\subsection{Salinity, Pollution, and Soil Remediation}

AI techniques enable the prediction of salinity distribution, pollutant transport, and remediation efficiency under varying climatic and soil conditions.

Soil salinisation constitutes a globally significant and accelerating form of soil degradation, affecting approximately 20\% of irrigated agricultural land and directly reducing plant-available water through osmotic stress \citep{nachshon2018cropland}. AI-based models have demonstrated strong predictive capability for mapping and monitoring soil salinity distribution from electromagnetic induction (EMI) survey data, satellite-derived spectral indices (e.g., the Salinity Index derived from SWIR reflectance), and combinations of environmental covariates including groundwater depth, evaporation rate, and irrigation history. Random forest and gradient boosting models trained on multi-depth EMI transect data and surface reflectance imagery have achieved spatial prediction accuracy ($R^2 > 0.80$) for soil electrical conductivity at field scales, providing the high-resolution salinity maps required for targeted leaching irrigation scheduling and salt-tolerant crop variety selection \citep{Lemenkova202515}.

Pollutant transport and groundwater contamination risk prediction represent a further domain in which AI models provide significant analytical value. Recurrent neural network models trained on groundwater quality time-series---including nitrate, pesticide, and heavy metal concentration records---have been used to predict future contamination levels under different irrigation and fertilisation scenarios, enabling proactive management of non-point source agricultural pollution \citep{rahmati2019groundwater, Lemenkova202214, Lemenkova202233}. Convolutional neural network models applied to hyperspectral imagery acquired over contaminated agricultural sites have demonstrated the capacity to map spatial distributions of cadmium, lead, and arsenic concentrations in topsoil with accuracy approaching laboratory analysis, at a fraction of the cost and time of conventional sampling campaigns. 

Physics-informed ML models that couple advection-dispersion equations with neural network surrogate functions for spatially variable hydraulic conductivity and retardation coefficients offer a computationally efficient alternative to finite-element solute transport simulators for scenario analysis of remediation strategies. AI-assisted optimisation of pump-and-treat and phytoremediation interventions---using genetic algorithms or reinforcement learning to identify remediation designs that minimise cost and duration for a target contaminant concentration threshold---is an emerging application with strong practical relevance for soil and groundwater restoration in agriculturally impacted landscapes.

\subsection{Digital Soil Mapping and Land Use Planning}

GIS-integrated AI frameworks support precision agriculture, land suitability analysis, and sustainable land-use planning. Digital soil mapping (DSM) uses spatially referenced soil profile observations as training data for AI models that predict the continuous spatial distribution of soil properties across unsurveyed areas from environmental covariate layers derived from remote sensing, terrain analysis, climate datasets, and geological maps \citep{mcbratney2003digital}. The SCORPAN framework---postulating that soil properties are functions of parent material, organisms, relief, climate, time, and spatial position---guides covariate selection in DSM, with terrain derivatives (slope, curvature, topographic wetness index, valley depth), multitemporal vegetation index composites, and lithological maps constituting the most consistently informative predictor groups across global benchmarking studies.

Random forest and gradient boosting models have dominated DSM benchmarking exercises, demonstrating particularly strong performance for mapping soil organic carbon, clay content, pH, bulk density, and available water capacity at 100--250~m resolution across national and continental domains \citep{hengl2017soilgrids250m}. The global SoilGrids250m product, produced using RF models trained on over 150,000 geo-referenced soil profiles and a comprehensive covariate stack, has established an internationally recognised baseline for high-resolution global soil property mapping. Deep learning architectures incorporating spatial context through CNN operations on image-structured covariate stacks, or through graph neural network propagation across spatial adjacency graphs of mapping units, are increasingly competitive with ensemble tree methods, particularly for structurally complex properties such as soil water retention parameters and field-scale hydraulic conductivity.

Multi-task learning frameworks that simultaneously predict multiple correlated soil properties in a shared representation space---exploiting the known pedological relationships between texture, organic carbon, pH, and hydraulic properties---achieve substantially lower prediction errors than independently trained single-property models, particularly in data-sparse regions where joint training provides implicit regularisation. GIS-integrated AI-based land suitability analysis frameworks that combine DSM outputs with hydrological, climatic, and agronomic constraints in a spatially explicit decision-support framework are increasingly deployed by national agricultural agencies to guide irrigation scheme planning, land-use zoning, and soil conservation investment prioritisation.

\subsection{Crop Yield Prediction and Precision Agriculture}

AI-driven yield forecasting systems integrate soil moisture, weather, and crop-growth data to improve productivity and resource management. AI-based crop yield prediction models combine soil moisture status, meteorological forcing, satellite-derived vegetation indices (NDVI, EVI, LAI), and agronomic management records to generate spatially explicit, in-season yield estimates \citep{you2017deep}. CNN--LSTM hybrid architectures that simultaneously capture the spatial pattern of remotely sensed crop reflectance and its temporal evolution through phenological development have achieved county-level maize and soybean yield prediction accuracies comparable to official agronomic survey estimates, several months before harvest. Transformer models equipped with cross-attention between the spectral and temporal dimensions of multi-date satellite image time series represent the current state of the art for in-season yield forecasting, demonstrating robust performance across diverse crop types, climatic zones, and management systems.

The integration of root-zone soil moisture dynamics---derived from LSTM-based prediction models informed by in situ sensors and satellite retrievals---as a key predictive covariate substantially improves yield forecast accuracy, particularly in rainfed systems and during drought years when soil water availability is the primary limiting factor on productivity. Precision variable-rate application of irrigation, guided by within-field soil moisture variability maps generated by ensemble ML models, has been demonstrated to reduce total irrigation water inputs by 15--30\% relative to uniform application while maintaining or increasing field-average yield, by directing water to the highest-yield-potential zones during critical crop growth stages \citep{vij2020iot}. Spatially explicit AI-based yield gap analysis---comparing predicted attainable yield under optimised soil moisture management with actual yield under observed conditions---provides quantitative diagnostics of where and why water-limited productivity shortfalls occur, enabling targeted interventions in precision agriculture programmes.

\subsection{Climate Change Adaptation}

ML models support drought forecasting, climate-risk assessment, and resilient agricultural planning under changing climatic conditions. Climate change is projected to substantially alter the spatial and temporal distribution of precipitation, intensify drought frequency and severity, increase evapotranspiration demand, and shift agroclimatic zones across major agricultural regions \citep{ipcc2022ar6}. AI models trained on historical soil--water observations and climate reanalysis data provide computationally efficient surrogates for the complex and nonlinear interactions between changing atmospheric forcing and soil hydrological response, enabling the rapid screening of climate change impacts and adaptation strategies across large ensembles of climate model projections \citep{trnka2022soil}.

ML emulators of physically-based crop--water models---trained to reproduce the input-output behaviour of DSSAT, AquaCrop, or SWAT simulations across a wide range of climate, soil, and management scenarios---can evaluate thousands of adaptation scenario combinations (drought-tolerant varieties, deficit irrigation strategies, cover cropping, soil carbon sequestration) within minutes, compared with weeks required for explicit physically-based simulation \citep{Meng}. Ensemble neural network models that incorporate CMIP6 climate model outputs as forcing inputs have been used to project mid-century changes in root-zone soil moisture deficits, growing season irrigation water requirements, and drought-induced yield losses for major staple crops under Representative Concentration Pathway scenarios, providing quantitative inputs for national adaptation planning and investment prioritisation \citep{trnka2022soil}.

Transfer learning frameworks that adapt AI models pretrained on historical climate--soil--yield relationships to future climate conditions---retraining on synthetic datasets generated from climate model projections to correct systematic biases in future forcing distributions---represent a promising approach for extending the temporal validity of data-driven agricultural models to climate states outside the historical training range. Physics-guided deep learning models that embed thermodynamic constraints on evapotranspiration and water balance closure as hard constraints provide an additional safeguard against physically implausible extrapolation under novel climatic conditions, ensuring that AI-generated climate adaptation assessments remain grounded in fundamental hydrological principles \citep{willard2022integrating}.

\section{\textcolor{blue}{Comparative Performance of AI Models}}

\subsection{Evaluation Metrics}

Rigorous comparative evaluation of AI models in soil--water applications requires a coherent set of quantitative performance metrics that capture different dimensions of predictive quality: overall error magnitude, bias, variance, and the capacity to reproduce the dynamic range and temporal structure of target variables. The selection of appropriate metrics is not merely a technical formality but fundamentally conditions which model properties are rewarded during training and how inter-model comparisons should be interpreted \citep{moges2021review}.

The Root Mean Square Error (RMSE) is the most widely reported accuracy metric in soil moisture prediction, evapotranspiration estimation, and streamflow forecasting applications. It measures the standard deviation of residuals between predicted ($\hat{y}_i$) and observed ($y_i$) values over $n$ samples:

\begin{equation}
RMSE=\sqrt{\frac{1}{n}\sum_{i=1}^{n}(y_i-\hat{y}_i)^2}
\end{equation}

Because RMSE weights large errors quadratically, it is particularly sensitive to outlier predictions and to errors during extreme events (heavy rainfall, severe drought), making it informative for applications where the tails of the error distribution are consequential \citep{moges2021review}. The Mean Absolute Error (MAE), defined as $MAE = n^{-1}\sum|y_i - \hat{y}_i|$, provides a complementary measure that treats all errors equally and is more robust to outliers, and is preferred when characterising typical prediction performance across a broad range of conditions.

The Nash--Sutcliffe Efficiency (NSE), originally proposed for hydrological model evaluation \citep{nash1970river}, normalises the mean squared error of the model against the variance of the observations:

\begin{equation}
NSE = 1 - \frac{\sum_{i=1}^{n}(y_i - \hat{y}_i)^2}{\sum_{i=1}^{n}(y_i - \bar{y})^2}
\end{equation}

where $\bar{y}$ is the observed mean. NSE values range from $-\infty$ to 1, with NSE~$= 1$ indicating perfect prediction and NSE~$= 0$ indicating that the model performs no better than the mean of the observations as a predictor. NSE is widely adopted in hydrological and soil--water modelling because it penalises both systematic bias and dynamic range errors, providing a single scalar summary of overall model fidelity. The Kling--Gupta Efficiency (KGE), a decomposition of NSE into correlation, bias ratio, and variability ratio components \citep{gupta2009decomposition}, is increasingly preferred in recent benchmarking studies because it provides a more diagnostically informative characterisation of model error structure and avoids certain known deficiencies of NSE in bias-dominated error regimes.

The coefficient of determination ($R^2$) measures the proportion of variance in the observations explained by the model and is ubiquitous in soil property prediction and digital soil mapping benchmarks. However, $R^2$ is insensitive to systematic multiplicative or additive bias---a model that consistently over- or under-predicts by a fixed factor can achieve $R^2 = 1.0$ while exhibiting substantial absolute error---and should therefore always be reported alongside RMSE or MAE. For probabilistic prediction outputs, the Continuous Ranked Probability Score (CRPS) and coverage probability of prediction intervals at specified nominal confidence levels provide complementary measures of uncertainty calibration that are essential for evaluating Bayesian and ensemble AI models applied to soil moisture forecasting \citep{gneiting2007strictly}.

\subsection{Comparative Analysis of Model Performance}
Recent studies indicate that deep learning architectures generally outperform traditional statistical and empirical models in highly nonlinear soil-water systems, Table \ref{tab:performance}. The mathematical variations and empirical performance boundaries of these predictive frameworks are systematized below. \textcolor{blue}{Of particular importance for soil moisture modelling is the class of physically constrained machine learning models, which embed physical laws such as the Richards equation, water balance constraints, and soil hydraulic relationships directly into the model training process. Recent studies have demonstrated that such physics-constrained approaches consistently outperform purely data-driven counterparts in predictive accuracy, physical consistency, and generalisation to sparse or extreme-value regimes \citep{ghorbani2021physics, Wang2023WR035194, Alsumaiei70551}.}

\textcolor{blue}{Synthesising quantitative performance indicators reported across the reviewed literature permits a more objective cross-model comparison than qualitative description alone. For soil moisture prediction at the point scale, LSTM and BiLSTM networks consistently achieve RMSE values of 0.02--0.05~m$^3$~m$^{-3}$ and $R^2$ of 0.85--0.97 in temperate and semi-arid climates, outperforming calibrated HYDRUS-1D simulations by 15--30\% in RMSE on the same station records \citep{kratzert2018rainfall}. CNN--LSTM hybrid models applied to field-scale spatiotemporal soil moisture mapping from Sentinel-1/2 inputs report $R^2$ of 0.82--0.93 and RMSE of 0.03--0.06~m$^3$~m$^{-3}$ under cross-validated evaluation. For evapotranspiration estimation, MLP and RF models achieve MAE values of 0.15--0.35~mm~day$^{-1}$ and NSE of 0.88--0.96 on FAO-56 reference ET benchmarks, while transformer-based architectures have been reported to attain NSE $> 0.95$ on multi-site streamflow and ET prediction tasks. For digital soil mapping, RF and XGBoost models predict soil organic carbon with $R^2$ of 0.60--0.80 and RMSE of 3--8~g~kg$^{-1}$ at 250~m resolution in continental benchmarks \citep{hengl2017soilgrids250m}, with deep learning CNN architectures achieving comparable or superior accuracy when spatially structured covariates are available. These quantitative benchmarks should be interpreted with caution, as they reflect heterogeneous experimental designs, geographic regions, soil types, and validation strategies; direct cross-study comparison is possible only when identical benchmark datasets and evaluation protocols are applied.}

%============= TABLE ==================================>
\footnotesize
\setlength{\LTleft}{0pt}
\setlength{\LTright}{0pt}

\begin{longtable}{@{\extracolsep{\fill}} p{0.15\textwidth} p{0.27 \textwidth} p{0.27\textwidth} p{0.13\textwidth} p{0.08\textwidth} @{}}
\caption{Comparative performance of AI models in soil--water applications. \label{tab:performance}}\\
\toprule
\textbf{AI Model} & \textbf{Key Advantages} & \textbf{Primary Limitations} & \textbf{Typical Accuracy} & \textbf{Key Citations} \\
\midrule
\endhead 

\bottomrule
\endfoot 

\textbf{Extreme Gradient Boosting (XGBoost)} & 
Exceptional handling of heterogeneous multi-source datasets (soil organic carbon, terrain indices) across multi-depth soil layers. & 
Prone to overfitting if hyperparameters (learning rate, max depth) are poorly optimized; high memory consumption. & 
$R^2$: 0.66--0.74\newline $NSE$: High & 
\citep{Demir2026} \\ \hline

\textbf{Temporal Convolutional Networks (TCN)} & 
Outperforms sequential models like LSTM by retaining longer effective history sizes using causal, dilated convolutions; highly parallelizable. & 
Demands a strictly regular, gap-free temporal sampling interval; can exhibit high inference latency during recursive prediction loops. & 
$R^2$: 0.96--0.99\newline $MAE$: $\approx$ 0.05 & 
\citep{ScientificReports2025} \\ \hline

\textbf{Ensemble Learning (RF + XGBoost Hybrid)} & 
Captures complex, highly non-linear interactions between crop leaf area index (LAI) and climate; improves $NSE$ metrics by up to 25\%. & 
Increased structural complexity leads to limited model interpretability and higher deployment friction at the IoT edge. & 
$R^2$: 0.81--0.84\newline $NSE$: $\ge$ 0.75 & 
\\ \hline

\textbf{Multilayer Perceptron (MLP)} & 
Demonstrates superior out-of-sample generalization capabilities for daily reference evapotranspiration ($ET_0$) in humid zones. & 
Extremely sensitive to input data scarcity; easily trapped in local minima during gradient descent backpropagation. & 
$R^2$: 0.95--0.98\newline $RMSE$: $\approx$ 0.21 & 
\citep{Acharki2025, Veredas2026} \\ \hline

\textbf{Support Vector Machines (SVM)} & 
Highly accurate for localized digital soil mapping of saturated soil hydraulic conductivity ($K_{sat}$) in complex headwater terrains. & 
Computationally expensive when scaled to vast, grid-based continental geographic data frameworks; poor text/symbolic processing. & 
$R^2$: 0.70--0.81\newline $RMSE$: Variable & 
\citep{MDPIReview2025} \\

\end{longtable}
\normalsize
%============= TABLE ==================================<

\subsection{Generalization and Transferability}

Transfer learning and domain adaptation techniques improve the scalability and robustness of AI models across diverse agricultural environments.

The generalisability of AI models beyond the specific geographic region, soil type, climate zone, and crop system represented in their training data is one of the most practically consequential challenges in agricultural AI \citep{reichstein2019deep}. A model that achieves high accuracy in cross-validated evaluation within its training domain but fails when applied to unseen conditions provides limited scientific and operational value, yet this failure mode is commonplace in the soil--water modelling literature where evaluation is frequently confined to the data-rich regions and environmental conditions that dominate the training corpus.

Transfer learning provides a principled approach to improving cross-domain model performance by leveraging knowledge acquired from a source domain (where abundant labelled training data are available) to improve learning in a target domain (where data are sparse). In the context of soil moisture prediction, models pretrained on dense networks of well-instrumented agricultural sites in North America or Central Europe can be fine-tuned on small amounts of locally collected data from target sites in data-sparse regions (sub-Saharan Africa, Central Asia), achieving substantially lower prediction errors than models trained from scratch on the limited target-domain data alone. The degree of transferability depends critically on the distributional similarity between source and target domains: when soil types, climate regimes, and vegetation systems differ markedly, domain adversarial training---which explicitly minimises feature distribution discrepancy between source and target domains---is necessary to prevent negative transfer, in which pretraining on a dissimilar source domain degrades target-domain performance \citep{pan2010survey}.

Large pretrained foundation models for earth observation---trained on petabyte-scale archives of multi-temporal, multi-sensor satellite imagery spanning diverse global land cover conditions---exhibit strong zero-shot and few-shot generalisation capacity for downstream soil and water property prediction tasks, owing to the breadth of spatiotemporal patterns encoded in their learned representations. Models such as the NASA-IBM Prithvi geospatial foundation model have demonstrated competitive fine-tuning performance on soil moisture downscaling, flood extent mapping, and crop type classification with substantially less task-specific training data than conventional architectures, indicating that large-scale pretraining effectively compresses transferable environmental knowledge that supports rapid adaptation to new domains \citep{jakubik2023foundation}. Developing rigorous evaluation protocols for transferability assessment---including spatially blocked cross-validation, leave-region-out evaluation, and prospective temporal holdout---is essential for honest reporting of model generalisability and for identifying failure modes before operational deployment.


\subsection{\textcolor{blue}{Decision Framework for Model Selection}}

\textcolor{blue}{A persistent gap in the AI soil--water modelling literature is the absence of structured guidance for practitioners on which modelling approach is most appropriate under a given combination of agricultural context, data availability, computational resource, and operational requirement. The following synthesis addresses this gap by mapping key selection criteria across the model families reviewed in this paper.}

\textcolor{blue}{When the primary constraint is training data scarcity (fewer than approximately 500 labelled observations, common in field-scale soil profile surveys), SVMs and shallow RF models should be preferred over deep learning architectures, as their regularisation mechanisms provide robust performance under small-sample conditions and avoid overfitting of high-capacity neural networks. When datasets are moderately sized (thousands of observations) and involve spatially structured covariates (remote sensing bands, terrain derivatives), RF and gradient boosting represent the most reliable choice: they handle mixed-type predictors and missing values natively, provide interpretable SHAP-based feature importance, and generalise well across diverse soil--climatic conditions. When the task involves long time-series with pronounced temporal dependencies---soil moisture forecasting, streamflow prediction, seasonal irrigation demand---LSTM and BiLSTM architectures are most appropriate, as their gated memory mechanisms model sequential dependencies over hundreds of time steps. For spatiotemporal field-scale tasks with satellite image inputs, CNN--LSTM hybrids or ConvLSTM architectures that simultaneously exploit spatial and temporal patterns are preferred. For multi-farm or privacy-sensitive deployments where raw data cannot be centralised, federated learning frameworks with FedAvg or personalised FL aggregation are the technically and legally appropriate architecture regardless of the base model family. When physical consistency of predictions is required---for regulatory compliance, hydrological balance closure, or deployment under extrapolative climate conditions---physics-informed neural networks or hybrid PINN--data-driven architectures should be selected even at the cost of additional implementation complexity. Interpretability requirements further condition model selection: if practitioners require human-interpretable explanations at the individual prediction level, intrinsically interpretable models (RF with SHAP, GAMs, decision trees) or post-hoc XAI wrappers applied to deep learning outputs should be incorporated from the design stage rather than as an afterthought.}

%%%%%%%%%%%%%%%%%%%%%%%%%%%%%%%%%%%%%%%%%%
%% SECTION 7: EXPLAINABLE AI, UNCERTAINTY,
%% AND ETHICAL CONSIDERATIONS
%%%%%%%%%%%%%%%%%%%%%%%%%%%%%%%%%%%%%%%%%%

\section{\textcolor{blue}{Explainable AI, Uncertainty, and Ethical Considerations}}

The deployment of AI and deep learning systems in agricultural water management raises important questions that extend beyond predictive accuracy to encompass interpretability, trustworthiness, and the equitable distribution of benefits and risks. Farmers, water resource managers, and policymakers who rely on AI-based decision-support tools require not only accurate predictions but also understandable explanations of why a model produced a particular output, honest quantification of the uncertainty surrounding that output, and assurance that the underlying system was designed, trained, and deployed in a manner consistent with principles of fairness, transparency, and accountability \citep{adadi2018peeking}. This section reviews the current state of explainable AI (XAI), uncertainty quantification, and ethical frameworks as they apply to AI-based soil--water systems.

% ----------------------------------------------------------
\subsection{Explainable Artificial Intelligence}
% ----------------------------------------------------------

The opacity of complex ML models---particularly deep neural networks with millions of parameters---has been a persistent concern since their widespread adoption. Although these models achieve impressive predictive accuracy, their internal representations are not directly interpretable by human experts, limiting their acceptance in high-stakes decision contexts such as irrigation management, drought early warning, and land-use planning, where the consequences of erroneous decisions can be severe and difficult to reverse. The field of explainable AI (XAI) has developed a rich toolkit of post-hoc and intrinsic explanation methods designed to make model behaviour interpretable without sacrificing predictive performance, Figure \ref{fig08}.

Among post-hoc explanation methods, SHapley Additive exPlanations (SHAP) \citep{lundberg2017unified} have emerged as the most widely adopted approach in soil and water science applications, owing to their rigorous game-theoretic foundation (based on Shapley values from cooperative game theory), their consistency across model classes, and their availability in efficient tree-based and kernel-based implementations. SHAP decomposes the prediction for each individual sample into additive contributions from each input feature, providing both local (sample-level) and global (dataset-level) attribution scores. 

In digital soil mapping applications, SHAP analysis has revealed that environmental covariates such as normalised difference vegetation index (NDVI), terrain curvature, and land surface temperature consistently rank among the most important predictors of soil moisture and organic carbon content across diverse geographic regions, providing physically interpretable validation of model behaviour \citep{padarian2019using}. 

In irrigation scheduling applications, SHAP-based attribution has been used to identify which meteorological variables---evapotranspiration demand, antecedent rainfall, vapour pressure deficit---most strongly drive irrigation timing recommendations, enabling agronomists to audit model decisions against domain knowledge \citep{zhou2023xai}.

\begin{figure}[H]
    \begin{center}
	\hspace{0pt}\includegraphics[width=14.0cm]{Figure_08.png}
    \end{center}
    \caption{Explainable AI framework for agricultural decision support. \textcolor{blue}{Software used: Inkscape, v. 1.4.4 for plotting data dashboard and clipart; draw.io for plotting taxonomy flowchart; GIMP v. 3.2.4 for plotting farm scene; PlotNeuralNet, NetworkX and Matplotlib for plotting DL architecture diagram. Noun Project for plotting conceptual OS architecture. Diagram} source: author.}
    \label{fig08}
\end{figure}

Local Interpretable Model-Agnostic Explanations (LIME) \citep{ribeiro2016lime} provide an alternative post-hoc approach in which a locally faithful linear surrogate model is fitted to the predictions of a complex model in the neighbourhood of each input sample. LIME is particularly useful for explaining predictions in image-based applications---for example, identifying which regions of a multispectral or SAR image most strongly influenced a soil classification or moisture retrieval prediction---and has been applied to CNN-based crop stress detection and soil type mapping from UAV and satellite imagery. Gradient-based saliency methods, including Gradient-weighted Class Activation Mapping (Grad-CAM) \citep{selvaraju2017gradcam} and integrated gradients, provide pixel-level attribution maps for CNN models, visually highlighting the image regions that most influenced a prediction.

Intrinsically interpretable model architectures---including decision trees, rule-based systems, and generalised additive models (GAMs)---remain valuable in soil--water applications where domain experts require full transparency of the decision logic. The Neural Additive Model (NAM) \citep{agarwal2021neural} extends GAMs to neural network feature functions while preserving additive interpretability, and has been applied to soil hydraulic property estimation with competitive accuracy relative to black-box deep learning models. Concept-based explanations, in which model predictions are expressed in terms of human-interpretable domain concepts (soil textural classes, drought severity indices, phenological stages), represent an emerging frontier that promises to bridge the gap between statistical model outputs and the conceptual vocabulary of agricultural practitioners.

Recent work has emphasised the importance of evaluating the faithfulness and robustness of XAI explanations: it has been shown that some popular attribution methods can produce misleading or unstable explanations when applied to deep networks with correlated input features---a common situation in soil--water datasets where meteorological and remote sensing variables exhibit strong multicollinearity \citep{adebayo2018sanity}. Developing XAI methods that are provably faithful to model behaviour under input correlation and that degrade gracefully under distribution shift remains an active and important research agenda for AI-based agricultural decision support.

% ----------------------------------------------------------
\subsection{Uncertainty Quantification}
% ----------------------------------------------------------

Quantifying the uncertainty of AI model predictions is as important as maximising their expected accuracy, particularly in risk-sensitive agricultural decision contexts where the tails of the predictive distribution determine the probability of crop failure, over-irrigation, or groundwater depletion. Predictive uncertainty in ML models arises from two distinct sources: \emph{aleatoric uncertainty}, which reflects irreducible noise in the measurement process (sensor error, natural variability) and cannot be reduced by collecting more training data; and \emph{epistemic uncertainty}, which reflects incomplete knowledge of the model parameters due to limited training data and can in principle be reduced with additional observations \citep{hullermeier2021aleatoric}.

Bayesian deep learning provides a principled framework for quantifying epistemic uncertainty by treating model weights as random variables with prior distributions that are updated to posterior distributions upon observing training data. In practice, exact Bayesian inference is computationally intractable for large neural networks, and approximate methods---including Monte Carlo Dropout \citep{gal2016dropout}, variational inference with the local reparametrisation trick, and Stochastic Weight Averaging Gaussian (SWAG)---are used to generate samples from an approximate weight posterior. When applied to soil moisture prediction and streamflow forecasting, Bayesian LSTM models with MC Dropout have been shown to produce well-calibrated predictive intervals that accurately capture the frequency of observed values falling within stated confidence bounds across diverse climatic conditions.

Deep ensemble methods \citep{lakshminarayanan2017simple}---in which a fixed number (typically 5--20) of independently initialised and trained models are combined by averaging their predictions---have demonstrated strong empirical performance as a practical alternative to Bayesian approximations, consistently producing sharper and better-calibrated uncertainty estimates than single models across multiple soil--water benchmarking studies. Conformal prediction \citep{angelopoulos2023conformal} offers a distribution-free, statistically rigorous framework for constructing prediction intervals with guaranteed marginal coverage under exchangeability assumptions, and is increasingly applied to spatiotemporal agricultural forecasting tasks where the assumption of independent and identically distributed (IID) test data may not hold. Communicating uncertainty to end users in agriculturally meaningful terms---for example, expressing soil moisture forecast uncertainty as the probability of crossing the wilting point threshold---remains an important and underexplored research challenge at the interface of AI, decision science, and agricultural extension \citep{moges2021review}.

% ----------------------------------------------------------
\subsection{Ethical and Socioeconomic Considerations}
% ----------------------------------------------------------

The deployment of AI-based soil--water management systems at scale raises profound ethical and socioeconomic questions concerning data ownership and sovereignty, algorithmic fairness, digital access inequality, and the distribution of risks and benefits across different categories of agricultural stakeholders \citep{janssen2020responsible, shepherd2022ethics}. These concerns are not peripheral to the technical enterprise of developing AI for precision agriculture but are central to determining whether the technology ultimately contributes to or undermines sustainable and equitable food systems.

Data ownership and privacy are among the most contentious issues in agricultural AI. Farm-level sensor data, yield maps, and soil survey results encode commercially sensitive information about land productivity and management practices. When this data is aggregated by commercial AI platform providers, farmers may inadvertently transfer proprietary agronomic intelligence to entities that could use it to inform commodity trading, land pricing, or the development of competing agricultural services \citep{WolfButtel}. Federated learning and differential privacy mechanisms \citep{dwork2014algorithmic} offer partial technical remedies by enabling AI model training without centralised data collection, but their adoption requires standardised data governance frameworks and legally enforceable data-sharing agreements that remain underdeveloped in most jurisdictions. The European Union's proposed Data Act and the FAIR (Findable, Accessible, Interoperable, Reusable) data principles \citep{wilkinson2016fair} provide important policy foundations, but their practical implementation in the agricultural sector is nascent.

Algorithmic fairness is a critical concern in AI systems that inform resource allocation decisions with distributional consequences. Soil--water AI models trained predominantly on data from technologically advanced, well-instrumented agricultural regions---as most current models are---systematically underperform in data-sparse smallholder farming contexts in sub-Saharan Africa, South Asia, and Latin America, where the consequences of irrigation mismanagement are most severe \citep{Xuetal2023}. This performance gap constitutes a form of algorithmic inequity that reinforces existing disparities in agricultural productivity and climate resilience. Addressing it requires deliberate investment in data collection, model adaptation, and co-design of AI tools with smallholder farming communities, respecting principles of participatory and responsible innovation \citep{janssen2020responsible}.

Digital infrastructure inequality---in connectivity, electrification, device availability, and data literacy---poses structural barriers to the adoption of AI-based precision irrigation in rural and low-income agricultural contexts. Designing AI systems that are robust to intermittent connectivity, operable on low-cost hardware, and explainable without specialist technical training is essential for ensuring that the benefits of agricultural AI are broadly accessible rather than accruing exclusively to large-scale, technologically endowed farm operators \citep{shepherd2022ethics}. The integration of local and indigenous soil and water knowledge into the design and validation of AI systems, alongside formal scientific datasets, represents both an ethical imperative and a scientifically productive strategy for developing models that are contextually appropriate and trusted by the communities they are intended to serve.

\section{\textcolor{blue}{Challenges and Limitations}}

\subsection{Data Scarcity and Quality}

Many agricultural regions lack high-quality datasets for robust AI model training and validation. The performance of AI and deep learning models is fundamentally bounded by the quality, volume, and spatial representativeness of their training data. In soil--water science, this dependency is particularly consequential because target variables---soil moisture, hydraulic conductivity, solute concentration, evapotranspiration---are continuous in space and time yet sampled at discrete, geographically biased locations. Existing soil profile databases are disproportionately concentrated in North America, Western Europe, and Australia; vast agricultural areas of sub-Saharan Africa, South Asia, and Central Asia are severely underrepresented \citep{hengl2017soilgrids250m}. ML models trained on spatially biased datasets exhibit systematically degraded prediction accuracy in underrepresented regions, propagating inequities in agricultural productivity data into global soil and water assessments.

Data quality limitations compound the scarcity problem. Sensor drift, cable damage, and power outages in field-deployed IoT networks introduce systematic gaps and outliers that corrupt training signals if not rigorously screened by automated quality-control algorithms \citep{gruber2020}. Label noise arising from measurement errors in soil reference samples, inconsistent analytical protocols between laboratories, and spatial support mismatch between point observations and gridded model targets imposes a practical accuracy ceiling that cannot be overcome by model architecture improvements alone. Temporal gaps caused by cloud cover in optical satellite time series are frequently non-random---cloud cover co-varies with rainfall events that are of primary hydrological interest---introducing systematic biases in models trained on cloud-filtered archives. Addressing these data infrastructure limitations through investment in globally distributed sensor networks, standardised measurement protocols, and open data sharing platforms is a prerequisite for realising the full predictive potential of AI in agricultural water management.

\subsection{Computational Complexity}

Large-scale distributed deep learning models require substantial computational resources and infrastructure. Training state-of-the-art deep learning architectures for soil--water applications---particularly transformer-based foundation models and spatiotemporal ensemble systems---demands computational resources that are unavailable in most agricultural research institutions and entirely inaccessible in developing-country contexts. The energy consumption of large-scale GPU cluster training carries a significant carbon footprint that is increasingly scrutinised from a sustainability perspective: estimates suggest that training a large transformer language model can emit hundreds of tonnes of CO$_2$ equivalent, a concern that extends to geoscience foundation models trained on comparable data volumes \citep{schwartz2020green}. Memory requirements for processing high-resolution spatiotemporal image sequences---as required for field-scale soil moisture downscaling or regional DSM---often exceed the capacity of individual GPU devices, necessitating model and data parallelism strategies with associated communication overhead and engineering complexity.

Inference latency introduces a further practical constraint for real-time irrigation control applications: edge-deployed AI processors on battery-powered IoT nodes must produce soil moisture estimates and irrigation decisions within milliseconds while consuming milliwatts of power, imposing severe model size and computational budget constraints that conflict with the accuracy requirements of precision management. Bridging the gap between the computational demands of high-accuracy deep learning models and the resource constraints of edge deployment is a fundamental systems engineering challenge that requires advances in hardware-aware neural architecture design, model compression, and energy-efficient inference acceleration.

\subsection{Model Interpretability}

Black-box AI systems often limit stakeholder trust and operational transparency. The opacity of deep neural network models constitutes a structural barrier to their adoption in agricultural decision-support contexts, where farmers, irrigation managers, and water regulatory agencies require not only accurate predictions but also intelligible explanations of the reasoning underlying model outputs \citep{arrieta2020explainable}. When an LSTM-based irrigation controller recommends withholding irrigation during a phenologically sensitive crop growth stage, or a gradient boosting model predicts a substantially elevated salinity risk in a specific field zone, practitioners need to understand which input factors drove these recommendations, whether the model is functioning within its domain of competence, and what the uncertainty bounds on its predictions are. Without this understanding, adoption of AI tools remains limited to technically sophisticated early adopters and fails to penetrate the broad agricultural practitioner community.

Current XAI methods---including SHAP, LIME, and gradient-based attribution---provide useful but imperfect approximations to model explanation: they can be computationally expensive, sensitive to correlated input features, and may produce misleading attributions under distribution shift \citep{adebayo2018sanity}. Developing XAI approaches that are provably faithful to model behaviour, efficient enough for routine use in operational agricultural AI systems, and expressible in terms of domain concepts (soil classes, phenological stages, drought severity indices) meaningful to agronomists and farmers remains an active and insufficiently resourced research frontier.

\subsection{Integration with Physical Models}

Hybrid frameworks integrating physics-based hydrological models with AI approaches remain an important research frontier. Despite the empirical success of purely data-driven AI models in soil--water prediction tasks, their inability to guarantee physically consistent outputs under distributional shift---predicting soil moisture values outside the physically plausible range, violating water balance closure, or producing conductivity estimates inconsistent with the van~Genuchten--Mualem framework---limits their reliability in operational hydrological applications and under novel climate forcing conditions \citep{willard2022integrating}. Physics-informed neural networks (PINNs) that embed governing partial differential equations as loss function constraints represent a technically promising but computationally challenging approach to addressing this limitation: optimising PINN loss functions that combine data fit and differential equation residual terms is numerically ill-conditioned for stiff PDEs such as the Richards equation, and practical PINNs for vadose zone flow remain an active research area. Modular hybrid frameworks---in which AI surrogate models replace computationally expensive components of established physically-based simulators while the overall water balance and process structure are maintained---offer a pragmatic near-term pathway for combining the physical consistency of process models with the data assimilation flexibility and pattern-recognition capacity of deep learning.

\subsection{Real-Time Deployment Constraints}

Latency, connectivity limitations, and energy requirements constrain real-time deployment in rural agricultural environments. The translation of AI-based soil--water models from research prototypes to operational field deployments is frequently impeded by infrastructure and connectivity constraints that are especially acute in smallholder and developing-country agricultural contexts. Reliable internet connectivity---required for cloud-based AI inference and remote model updating---is unavailable across large fractions of the world's agricultural land, particularly in sub-Saharan Africa and South Asia where water management improvements would generate the greatest food security benefits \citep{shepherd2022ethics}. Edge AI deployment on locally powered IoT nodes mitigates this connectivity dependency but introduces its own constraints: model inference must complete within strict latency budgets; hardware must withstand field temperature extremes, humidity, and mechanical vibration; and power must be supplied by solar panels or batteries with limited capacity. Firmware updates for edge-deployed models across spatially distributed sensor networks require secure over-the-air update protocols and robust version management systems that represent a significant operational engineering burden. Harmonising these deployment realities with the accuracy and latency requirements of precision irrigation control constitutes a systems-level challenge that demands closer collaboration between AI researchers and agricultural engineering practitioners.

\subsection{\textcolor{blue}{Model Validation, Uncertainty Analysis, and Data Quality Effects}}

\textcolor{blue}{The practical adoption of AI-based soil--water systems depends critically on rigorous model validation procedures and transparent uncertainty quantification, yet these topics receive insufficient coverage in most published studies. A persistent methodological weakness in the reviewed literature is the use of random cross-validation splits that do not account for spatial autocorrelation in soil property datasets: when training and test samples are spatially proximate, cross-validated accuracy estimates are optimistically biased because the model can exploit spatial structure rather than learning generalisable pedological relationships. Spatially blocked cross-validation---partitioning observations into spatially contiguous folds that enforce minimum separation distances between training and test sets---is the methodologically correct evaluation strategy for soil mapping and field-scale prediction tasks and should be adopted as the standard reporting requirement. Temporal holdout validation, in which models are evaluated on complete seasons or years withheld from training, is the appropriate strategy for time-series soil moisture and irrigation forecasting tasks, as it tests the model's ability to generalise forward in time under potentially novel forcing conditions.}

\textcolor{blue}{Sensitivity analysis of AI model predictions with respect to input data quality---quantifying how much prediction error increases as a function of sensor noise level, temporal gap frequency, spatial sampling density, or label measurement error---provides essential guidance for field deployment design and data collection investment decisions, yet is rarely reported alongside accuracy benchmarks. Uncertainty quantification frameworks---whether Bayesian (Monte Carlo Dropout, variational inference), ensemble-based (deep ensembles, random forest prediction intervals), or conformal prediction---should be incorporated into AI soil--water models from the design stage rather than added post-hoc, as the predictive uncertainty is as operationally consequential as the point estimate for irrigation decisions near the wilting point threshold or drought early-warning triggers. The effect of training data quality on AI model prediction reliability is further compounded by label noise: soil moisture measurements from field sensors carry calibration uncertainties of 0.02--0.05~m$^3$~m$^{-3}$, which sets a practical accuracy floor that cannot be overcome by model architecture improvements, and should be explicitly acknowledged when reporting RMSE values approaching this range. Developing standardised reporting guidelines for validation design, uncertainty communication, and sensitivity analysis in AI soil--water modelling studies would substantially improve the comparability and practical credibility of published results.}

\section{\textcolor{blue}{Future Research Directions and Technical Horizons}}

The systemic bottlenecks hindering widespread field implementation of agricultural AI frameworks, alongside corresponding next-generation research opportunities, are detailed below.

\textcolor{blue}{Despite rapid progress, the field faces several unresolved challenges that define specific and actionable future research priorities. First, the absence of standardised benchmark datasets for soil moisture prediction, pedotransfer function development, and digital soil mapping impedes rigorous cross-study comparison and slows the cumulative scientific progress achievable through reproducible benchmarking; establishing openly accessible, globally distributed benchmark corpora analogous to ImageNet for computer vision represents a critical infrastructure investment. Second, improving the transferability of AI models across agroecological regions requires systematic investigation of domain adaptation architectures---including meta-learning, domain adversarial training, and few-shot fine-tuning strategies---benchmarked against consistent leave-region-out evaluation protocols that honestly assess cross-regional generalisation failure modes. Third, the integration of physics-informed AI models with operational precision irrigation systems requires development of computationally efficient PINN architectures capable of training at field-to-catchment scales, with particular attention to the numerical stiffness of the Richards equation and the spatial heterogeneity of vadose zone hydraulic properties. Fourth, advancing privacy-preserving federated learning for agricultural decision support requires development of communication-efficient aggregation algorithms robust to the extreme non-IID data distributions arising from soil and climate heterogeneity across participating farm nodes, as well as legally enforceable data governance frameworks aligned with national agricultural data sovereignty policies.}

%============= TABLE ==================================>
\footnotesize
\setlength{\LTleft}{0pt}
\setlength{\LTright}{0pt}

\begin{longtable}{@{\extracolsep{\fill}} p{0.10\textwidth} p{0.32\textwidth} p{0.32\textwidth} p{0.10\textwidth} @{}}
\caption{Future research gaps and opportunities in AI-driven agriculture. \label{tab:future_gaps}}\\
\toprule
\textbf{Research Domain} & \textbf{Identified Knowledge \& Technical Gap} & \textbf{Emerging Opportunity \& Solution} & \textbf{Key Sources} \\
\midrule
\endhead 

\bottomrule
\endfoot 

\textbf{Model Interpretability} & 
Standard deep architectures function as ``black boxes,'' causing low practitioner trust when models output highly irregular irrigation schedules. & 
Development of Explainable AI (XAI) frameworks (e.g., SHAP, Integrated Gradients) to map physical features to neuron activations. & 
\citep{Pai_2025} \\ \hline

\textbf{Physical Consistency} & 
Purely data-driven models frequently output predictions that violate fundamental laws of thermodynamics and mass-conservation. & 
Integration of Physics-Informed Neural Networks (PINNs) that bake Darcy's law and Richards equation directly into the loss function. & 
 \\ \hline

\textbf{Data Privacy \& Sovereignty} & 
Centralized cloud storage of granular farm telemetry creates corporate surveillance vulnerabilities and data ownership disputes. & 
Implementing Federated Learning (FL) to train localized models on edge gateways without moving raw agricultural data off-site. & 
\citep{electronics14132512} \\ \hline

\textbf{Cross-Regional Generalization} & 
Models optimized for a specific soil type or micro-climate exhibit drastic performance degradation when deployed in unseen regions. & 
Development of multi-modal Agricultural Foundation Models pre-trained on continental remote sensing and fine-tuned locally. & 
\\ \hline

\textbf{Edge Compute Constraints} & 
High-fidelity transformer networks are computationally too heavy for low-power, solar-powered microcontroller valves in the field. & 
Application of hardware-aware Knowledge Distillation and structural network pruning to compress multi-million parameter models. & 
\\

\end{longtable}
\normalsize
%============= TABLE ==================================<

\subsection{Physics-Informed AI}

Physics-informed neural networks (PINNs) offer opportunities to integrate physical laws with data-driven learning for improved generalization. The integration of domain physical knowledge with deep learning through physics-informed neural networks (PINNs) and physics-guided loss functions represents one of the most promising frontiers for improving the reliability and generalisability of AI-based soil--water models \citep{karniadakis2021physics}. 

By embedding differential equation constraints---such as the Richards equation governing unsaturated flow, Darcy's law relating flux to hydraulic gradient, or the water balance equation linking precipitation, evapotranspiration, runoff, and soil moisture storage change---directly into the neural network loss function, PINNs constrain predictions to remain consistent with known physical laws, even in data-sparse regions where purely data-driven models are prone to physically implausible extrapolation.  The combination of PINNs with ensemble uncertainty quantification and Bayesian parameter inference frameworks would provide a principled mechanism for simultaneously learning soil hydraulic parameters, predicting soil moisture dynamics, and quantifying prediction uncertainty within a unified physics-consistent model. Hybrid architectures that couple PINN-based PDE solvers with data-driven parameterisation of spatially variable hydraulic properties---learning the spatially distributed van~Genuchten parameters as neural network outputs rather than prescribing them from PTFs---represent a particularly compelling direction for reducing the parameter estimation burden of physically-based soil--water models without sacrificing their physical interpretability.

\subsection{Federated Learning in Agriculture}

Federated learning can facilitate collaborative agricultural intelligence while preserving privacy and decentralization. Federated learning (FL) offers a technically and legally robust mechanism for aggregating AI model training across geographically distributed farms, research stations, and national agricultural networks without requiring sensitive farm-level data to be centralised or disclosed. The agricultural domain presents distinctive challenges for FL that motivate important future research directions. Non-IID data heterogeneity---arising from the pronounced variation in soil types, climate regimes, crop species, and management practices across participating nodes---causes significant performance degradation under standard FedAvg aggregation, necessitating personalised FL approaches that maintain local model components adapted to site-specific conditions while sharing globally beneficial representations. 

Communication-efficient FL algorithms that compress gradient updates through quantisation, sparsification, or local SGD with infrequent synchronisation are needed to accommodate the limited and intermittent bandwidth of rural agricultural networks. The combination of FL with differential privacy mechanisms---adding calibrated noise to local model updates before transmission to prevent inference of individual farm records from the aggregate model---requires careful calibration to balance privacy protection against the accuracy degradation imposed by noise injection, particularly for small-farm nodes with limited local training data. Future FL systems for agricultural AI should integrate standardised data harmonisation protocols, interoperable communication interfaces, and transparent audit mechanisms that enable participating farms to verify the integrity of the aggregation process and retain meaningful control over their data contributions, in alignment with emerging agricultural data governance frameworks.

\subsection{Edge AI and Smart Farming}

Edge computing and low-power AI devices will enable autonomous irrigation and real-time field-scale analytics. The continuing miniaturisation and energy efficiency improvements of neuromorphic and application-specific integrated circuit (ASIC) hardware, combined with advances in model compression through knowledge distillation, quantisation-aware training, and neural architecture search, are progressively closing the gap between the computational requirements of state-of-the-art deep learning models and the resource budgets of field-deployable edge processors. 

Future smart farming systems will likely comprise hierarchical edge-fog-cloud AI architectures in which ultra-low-power microcontrollers at the sensor node level run highly compressed models for real-time anomaly detection and actuation control, Figure \ref{fig09}.

\begin{figure}[H]
    \begin{center}
	\hspace{0pt}\includegraphics[width=14.0cm]{Figure_09.png}
    \end{center}
    \caption{Future ecosystem of AI-enabled sustainable agriculture. \textcolor{blue}{Software: Midjourney v. 6.1 for infographic illustration of farm scenes; Inkscape, v. 1.4.4 for figure composition and layout}. Diagram source: author.}
    \label{fig09}
\end{figure}

Among other, intermediate fog gateway processors aggregate sensor streams and run more complex soil moisture and evapotranspiration prediction models; and cloud or HPC resources periodically retrain global models and push updated lightweight models to the edge tier. Autonomous variable-rate irrigation systems---comprising distributed wireless sensor arrays, AI-driven soil moisture prediction at sub-field resolution, weather-forecast-integrated optimisation, and automated valve actuation---represent a realistically near-term technology convergence that could substantially reduce agricultural water consumption in irrigated systems. Standardising sensor communication protocols (LoRaWAN, NB-IoT, MQTT), data formats, and AI model interchange standards across hardware vendors and platform providers is essential to enable the interoperability required for scaling these systems from experimental deployments to commercial precision agriculture ecosystems.

\subsection{Digital Twins for Soil--Water Systems}

Digital twins (DT) can simulate soil-water interactions and optimize agricultural management under dynamic environmental conditions. DT technology---creating real-time, continuously updated virtual replicas of physical agricultural systems that synchronise observations from sensor networks and satellite platforms with AI and physically-based simulation models---offers transformative potential for soil--water management and climate adaptation planning \citep{pylianidis2021digital}. An agricultural digital twin integrates data streams from field sensors, satellite observations, weather forecasts, and agronomic management records into a dynamically updated state representation of the soil--water--plant--atmosphere system, which serves as both a predictive engine for forward simulation under different irrigation and climate scenarios and a decision-support interface for farm managers and water resource planners. 

ML surrogate models trained to emulate high-fidelity HYDRUS or SWAT simulations at low computational cost are the enabling technology for real-time DT updating: they provide sub-second response times necessary for interactive what-if scenario exploration while maintaining quantitative consistency with the physical process representations embedded in the reference simulators. Future research should address the development of DT architectures that quantify and propagate model and observational uncertainties throughout the virtual system state, provide rigorous validation frameworks for assessing DT predictive fidelity under novel forcing conditions, and support multi-stakeholder decision-making across the spatial scales relevant to both individual farm management and watershed-level water allocation governance. The integration of agricultural DTs with climate service platforms---providing dynamically downscaled seasonal climate forecasts as boundary conditions for soil--water system simulation---would enable adaptive water management that anticipates seasonal rainfall variability months in advance rather than merely reacting to observed soil moisture deficits.

\section{\textcolor{blue}{Conclusions}}
Artificial intelligence, ML, and distributed deep learning are transforming soil science and sustainable agricultural water management. AI-driven frameworks provide powerful tools for soil moisture prediction, irrigation optimization, digital soil mapping, and climate resilience. Despite significant progress, challenges related to data quality, interpretability, computational requirements, and real-time deployment remain unresolved. Future research should emphasize explainable AI, physics-informed learning, federated intelligence, and integrated digital agricultural ecosystems to support resilient and sustainable food production systems. Future research should prioritise the development of computationally efficient PINN architectures tailored to the numerical stiffness and spatial heterogeneity characteristic of vadose zone flow, including adaptive collocation point sampling strategies and domain decomposition methods that enable PINN training at field-to-catchment scales. 


%----------- section ----------->
\authorcontributions{Supervision, conceptualization, methodology, software, resources, funding acquisition, and project administration, writing---original draft preparation, methodology, software, data curation, visualization, formal analysis, validation, writing---review and editing, and investigation, P.L.}

%\funding{}

\institutionalreview{Not applicable.}

\informedconsent{Not applicable.}

\dataavailability{Not applicable} 

\acknowledgments{The authors thank the reviewers for reading and review of this manuscript.}

\conflictsofinterest{The authors declare no conflict of interest.} 

\abbreviations{Abbreviations}{
The following abbreviations are used in this manuscript:\\

\noindent 
\begin{tabular}{@{}ll}
AI   & Artificial Intelligence\\
ANN  & Artificial Neural Network\\
BiLSTM & Bidirectional Long Short-Term Memory\\
CNN  & Convolutional Neural Network\\
DL   & Deep Learning\\
DOAJ & Directory of Open Access Journals\\
DSM  & Digital Soil Mapping\\
DQN  & Deep Q-Network\\
FL   & Federated Learning\\
GNN  & Graph Neural Network\\
IoT  & Internet of Things\\
IPCC & Intergovernmental Panel on Climate Change\\
LSTM & Long Short-Term Memory\\
MAE  & Mean Absolute Error\\
MDPI & Multidisciplinary Digital Publishing Institute\\
ML   & \textcolor{blue}{Machine Learning}\\
MLP  & Multi-Layer Perceptron\\
MODFLOW & Modular Groundwater Flow Model\\
NSE  & Nash--Sutcliffe Efficiency\\
PAWC & Plant-Available Water Capacity\\
PINN & Physics-Informed Neural Network\\
PTF  & Pedotransfer Function\\
RF   & Random Forest\\
RL   & Reinforcement Learning\\
RMSE & Root Mean Square Error\\
RNN  & Recurrent Neural Network\\
SAR  & Synthetic Aperture Radar\\
SHAP & SHapley Additive exPlanations\\
SMAP & Soil Moisture Active Passive\\
SOC  & Soil Organic Carbon\\
SVM  & Support Vector Machine\\
SVR  & Support Vector Regression\\
SWAT & Soil and Water Assessment Tool\\
TDR  & Time-Domain Reflectometry\\
UAV  & Unmanned Aerial Vehicle\\
ViT  & Vision Transformer\\
XGBoost & Extreme Gradient Boosting\\
\end{tabular}
}

%%%%%%%%%%%%%%%%%%%%%%%%%%%%%%%%%%%%%%%%%%
\begin{adjustwidth}{-\extralength}{0cm}
%\printendnotes[custom] % Un-comment to print a list of endnotes

\reftitle{References}
%\nocite{*}

%=====================================
% References, variant A: external bibliography
%====================================

\bibliography{Bibliography.bib}

%%%%%%%%%%%%%%%%%%%%%%%%%%%%%%%%%%%%%%%%%%
\end{adjustwidth}
\end{document}
